\section{Finite-index reductions and certificate conditions}
\label{ms:finite:scope}
The finite-index arguments use one standardized iid marginal at each actual
index, with the notation $S_n=X_1+\cdots+X_n$ and
$R_n(P)=\sqrt n\,\Delta_n(P)/\rho(P)$. The notation $C_n(P)$ below
also denotes $R_n(P)$; $C_n$ is its supremum over standardized laws.
The numerical conclusions are established by the verified certificate conditions
in the public release (\texttt{full\_pass: "ESTABLISHED"}).

The first part develops the conditional-count bounds shared by several
finite-index comparisons. The small-index reductions then proceed by a
closed cubic dual at $n=1$ (Section~\ref{ms:finite:n1}), a rational tensor
majorant at $n=2$ (Section~\ref{ms:finite:n2:section}), independence-null
polynomials at $n=3$ (Section~\ref{ms:finite:n3:section}), a canonical
reference law and physical Prawitz comparison at $n=4$
(Section~\ref{ms:finite:n4all:section}), and a moment-domain cover with a
gap-free count comparison at $n=5$ (Section~\ref{ms:finite:five:section}).
The later sections supply signed moment and gap refinements, an exact phase
helper for $n\in\{6,7,8,9\}$, a B1 sine interface for $n=26,\ldots,30$,
and conditional resource transfers at floors $25$, $31$ and $40$.
Together with the complete thirteen-role partition in Section~\ref{ms:certificates}
(\texttt{N1}, \texttt{G2}, \texttt{G3}, \texttt{G4/ROOT4}, \texttt{C5-one-free},
\texttt{projected-C69}, \texttt{projected-C1019}, \texttt{current-C2024},
\texttt{current-C25}, \texttt{current-C2639}, \texttt{current-C4042},
\texttt{current-C4348}, and \texttt{independent-tail49}), these reductions and
verified certificates cover all finite indices and the unbounded tail.

In count calculations $m=\E X_+$ and $h$ is the reciprocal span.
In the later index-transfer section $m$ denotes the actual integer index,
and the first positive mean is written $\mu$. The symmetrization section
states its separate first-absolute-moment convention for $h$.

\begin{assumption}[Analytic inputs and their scope]
\label{ms:finite:inputs}
Each application requires the following inputs at its actual index, for the
same marginal, signed branch, raw threshold and moments.
\begin{enumerate}
\item\label{ms:finite:input:cut49:global:attainment}
\label{ms:finite:input:cut49:global:contact}
\label{ms:finite:input:cut49:global:finitecaps}
Global contrary-witness attainment, positive closed-edge orientation,
contamination contact, and the stated global caps, for the unrestricted
maximizing law.
\item\label{ms:finite:input:cut49:support:theorem}
The actual-index contact exclusion on both complete support rays and the
whole threshold line, with every Taylor remainder and tail included and
favourable coefficients evaluated on their stated index domains. The separate
two- and three-index support proofs are given below.
\item\label{ms:finite:input:orig:thm:defect}
The sharp moment-domain inequality
$\rho^2(\gamma^2+4)\ge(\gamma^2+2)^2$, together with its stated moment
hypotheses. Its use only enlarges the realizable moment domain.
\item\label{ms:finite:input:orig:lem:BPcardinal}
\label{ms:finite:input:afcore:eq:prawitz}
\label{ms:finite:input:afcore:sec:fourier}
The closed-CDF Prawitz majorant with the exact normalization used below,
and the complete four-term Fourier bound, including the infinite normal tail.
\item\label{ms:finite:input:afcore:eq:gap-definitions}
\label{ms:finite:input:afcore:prop:conserved}
\label{ms:finite:input:afcore:sec:conserved}
The gap functional with its half-open endpoint convention and conserved
contact identity for the same law, including all nonnegative slack terms
that are discarded in the displayed inequalities.
\item\label{ms:finite:input:rescore:re-enclosure}
\label{ms:finite:input:rescore:sec:setup}
\label{ms:finite:input:rescore:u-contract}
\label{ms:finite:input:rescore:u-three-uses}
\label{ms:finite:input:rescore:u-phase-interface}
The resource enclosure and finite propagation implications on their original
index domains, with a genuine acquired central gap, complete signed resource
fields, positive divisors and the same physical defect throughout.
\item\label{ms:finite:input:cut49:short:arithmetic}
\label{ms:finite:input:elementary:section}
Directed rational enclosures of elementary constants, Gaussian functions and
trigonometric functions, with analytic remainder bounds and their domain gates.
\end{enumerate}
\end{assumption}
The corresponding analytic arguments are located as follows. The attainment and
contact proofs are in Section~\ref{ms:extremizers}; the external moment caps
are in \eqref{ms:structural-classical}; and the sharp moment domain is proved
in Section~\ref{ms:moment-domain}. The conserved identity is in
Section~\ref{ms:conserved}. The Prawitz construction, Fourier bounds
and resource enclosure are in Sections~\ref{ms:fourier:kernel}
and~\ref{ms:fourier:afcore:sec:fourier}, including
Section~\ref{ms:fourier:rescore:re-enclosure}. Every application also retains the numerical conditions and index restrictions
of the result it uses.

\subsection{The conditional count comparison}
\label{ms:finite:countcore:sec}

We develop the outgoing and incoming probability bounds used after support
and resource acquisition, keeping the same marginal throughout.

\subsubsection{One marginal, all conditional counts}

Fix $n\ge2$ and a standardized marginal with $\rho<\infty$. Use the sign coordinates of
\eqref{ms:bernoulli-coordinates}: $p=\Pp(X>0)$, $q=1-p$, $v=pq$,
$m=\E X_+=\E X_->0$, with zero atoms assigned to the nonpositive
class. In the reciprocal-span notation,
\[
 h=\ell^{-1}=v/m=\sqrt{v+s},\qquad I=B_0.
\]
Thus the main representation becomes
\begin{equation}
 hX=I-p+U,\qquad \E(U\mid I)=0.
 \label{ms:finite:countcore:representation}
\end{equation}
Retain the same conditional laws $U_0,U_1$ and moments
$w_i=\E U_i^2$, $\mu_i=\E U_i^3$. Their total noise is
$s=qw_0+pw_1=h^2-v$ and $\lambda=ns$. If $s>0$, write
$\alpha=pw_1/s\in[0,1]$, so that
$w_0=s(1-\alpha)/q$ and $w_1=s\alpha/p$.
At $s=0$ both conditional laws vanish and any $\alpha\in[0,1]$
is allowed, without division by $s$. Higher moments below require the
stated bounded-support hypotheses.

Use the conditional sums $W_j$ from \eqref{ms:count-identity}, and write
$b_j(p)=\binom njp^jq^{n-j}$. Conditional independence gives
\begin{equation}
 V_j:=\E W_j^2=(n-j)w_0+jw_1
 =\lambda\left(1+\frac{(j-np)(\alpha-p)}{npq}\right).
 \label{ms:finite:countcore:variance}
\end{equation}

\begin{lemma}[Closed count comparison]
\label{ms:finite:countcore:countidentity}
For \(z=np+hH=k+u\), with \(0\le k\le n\),
\begin{equation}
 \Pp(S_n\le H)
 \le F_{\mathrm{Bin}(n,p)}(k)-b_k(p)\Pp(W_k>u)
       +\sum_{j>k}b_j(p)\Pp(W_j\le u-(j-k)).
 \label{ms:finite:countcore:countineq}
\end{equation}
The outgoing event is strict and every incoming event is closed.
\end{lemma}
\begin{proof}
In the exact signed identity \eqref{ms:count-difference}, discard the
nonpositive outgoing terms with count $j<k$ and retain the $j=k$ term.
The strict/closed conventions in that identity include every event-plane atom.
\end{proof}

All ensuing formulas hold on either probability half. On a closed interval
\([l,r]\subset(0,1)\), the exact range of \(v=p(1-p)\) is
\[
 v_- =\min\{l(1-l),r(1-r)\},\qquad
 v_+=\begin{cases}1/4,&l\le1/2\le r,\\
                  \max\{l(1-l),r(1-r)\},&\text{otherwise}.
       \end{cases}
\]
The maximum of \(b_j\) is attained at \(j/n\) if that point belongs
to the interval, and otherwise at an endpoint; its minimum is at an
endpoint. Also
\[
 \frac{d}{dp}F_{\mathrm{Bin}(n,p)}(k)
 =-\frac{(n-k)b_k(p)}{1-p}\le0.
\]
Thus no reflection of \(p\), the count, or the signed moments is needed
when a parameter interval meets or lies to the right of \(1/2\).

\subsubsection{Moment identities and necessary contractions}

Suppose the acquired conditional supports are
\begin{equation}
 U_0\in[-O_0,a_0],\qquad U_1\in[-a_1,O_1],
 \qquad O_i,a_i\ge0.
 \label{ms:finite:countcore:support}
\end{equation}
For a centered variable \(T\in[-L,R]\), of variance \(w\),
integration of nonnegative polynomials gives
\begin{equation}
 w\le LR,\qquad
 w^2/L-Lw\le\E T^3\le Rw-w^2/R,
 \qquad \E T^4\le(R-L)\E T^3+LRw.
 \label{ms:finite:countcore:momentbounds}
\end{equation}
For the cubic bounds use respectively
\((T+L)(T-w/L)^2\ge0\) and
\((R-T)(T+w/R)^2\ge0\); for the fourth bound multiply
\((T+L)(R-T)\ge0\) by \(T^2\). If either endpoint is zero,
centering forces \(T=0\); all corresponding moment terms vanish.

For later substitutions, the identities \eqref{ms:conditional-moments}
read as follows in reciprocal-span units:
\begin{align}
 \rho h^3&=v(1-2v)+3v(w_0+w_1)-q\mu_0+p\mu_1,
       \label{ms:finite:countcore:rhoidentity}\\
 \gamma h^3&=v(1-2p)+3v(w_1-w_0)+q\mu_0+p\mu_1,
       \label{ms:finite:countcore:gammaidentity}\\
 \beta&=\frac{v(1-2p)+s(2\alpha-1)}{v+s},\qquad
 m=\frac v{\sqrt{v+s}},
       \label{ms:finite:countcore:betaidentity}
\end{align}
Here $\gamma=\E X^3$ and $\beta=\E[X|X|]$; the differing signs
of $\mu_0$ follow from the two actual sign classes. The expressions for
$\beta$ and $m$ use $pw_1-qw_0=s(2\alpha-1)$ and $h^2=v+s$.

These identities give necessary closed-domain contractions. For example,
if \(m_-\le m\le m_+\), then
\[
 \lambda=n(v^2/m^2-v),\qquad v\ge m_-^2.
\]
On the resulting closed variance interval a lower bound is obtained by
minimizing \(n(v^2/m_+^2-v)\), and, when \(m_->0\), an upper bound
by maximizing \(n(v^2/m_-^2-v)\). These are convex quadratics: the
minimum occurs at an endpoint or its stationary point, and the maximum
at an endpoint. If \(\theta\) is the within-sign variance of \(X\), then
\[
 \theta=\frac{s}{v+s},\qquad
 \lambda=\frac{nv\theta}{1-\theta}.
\]
Thus \(\theta\le\Theta<1\) implies
\(\lambda\le nv_+\Theta/(1-\Theta)\).

If a simultaneous acquired resource gives
\begin{equation}
 \frac{a_0+a_1}{d_0}+\frac\lambda\Lambda\le1,
 \qquad d_0,\Lambda>0,
 \label{ms:finite:countcore:jointbudget}
\end{equation}
then \eqref{ms:finite:countcore:momentbounds} also implies
\[
 \lambda\le
 \left\{\frac1\Lambda+\frac1{nd_0}
 \min_{p,\alpha}\left[
 \frac{1-\alpha}{(1-p)O_0}+\frac\alpha{pO_1}\right]\right\}^{-1}.
\]
Here the minimum is over the incoming probability/allocation rectangle.
First minimize at an \(\alpha\) endpoint. For each such endpoint the
remaining function is \(A/(1-p)+B/p\), with \(A,B\ge0\); its
minimum is at an endpoint or the stationary value
\(A+B+2\sqrt{AB}\). A smaller rational lower bound for this minimum
remains valid. These independent caps intersect the working
\(\lambda\) interval; they do not replace the intercept \(\Lambda\)
in \eqref{ms:finite:countcore:jointbudget}.

Signed \(\beta\) bounds in \eqref{ms:finite:countcore:betaidentity} similarly
restrict \(\alpha\), or give a necessary positive \(\lambda\) floor.
For example, if \(\beta\le b<1-2p_+\), then
\[
 s(b+1-2\alpha)\ge v(1-2p-b)
 \ge v_-(1-2p_+-b)>0.
\]
A nonpositive upper bound for the left coefficient excludes the law;
otherwise division gives a valid lower bound for \(s\).
Every exclusion uses a strict contradiction. Equal interval endpoints
and zero variance remain included.

There is also a rho-derived noise floor. The rough cubic bounds
\(-\mu_0\le O_0w_0\), \(\mu_1\le O_1w_1\) in
\eqref{ms:finite:countcore:rhoidentity} imply
\[
 r(v+s)^{3/2}\le v(1-2v)+Ms,\qquad
 M=\max_{p,\alpha\ \mathrm{corners}}
       \{(3p+O_0)(1-\alpha)+(3q+O_1)\alpha\},
\]
when \(\rho\ge r\). The maximum is at rectangle corners because the
expression is affine in each variable. Convexity gives
\[
 A(v):=rv^{3/2}-v+2v^2
 \le (M-\tfrac32r\sqrt v)s.
\]
Let \(t\le\sqrt{v_-}\) be a nonnegative rational lower enclosure and
\(A_-=rv_-t-v_-+2v_-^2>0\). Then
\(rt>1-2v_-\), so
\(A'(v_-)\ge3rt/2-1+4v_->1/2+v_->0\).
Since \(A''>0\), \(A(v)\ge A_-\) on the whole variance interval.
Thus \(M-3rt/2\le0\) is impossible; otherwise
\(\lambda\ge nA_-/(M-3rt/2)\). If the positive \(A_-\) gate fails,
this optional floor is simply omitted.

\subsubsection{The physical-defect halfplane}

Assume compact support and a genuine central gap $(-A,B)$, with
$A,B>0$. Use $\tau,J,D$ from \eqref{ms:physical-defect-definition},
$\sigma=\sqrt{10}$, and the nonnegative residuals $T_\pm,R_\pm$ of
Lemma~\ref{ms:physical-defect}. Its outer-gap estimates and the
identity \eqref{ms:defect-decomposition} give
\begin{equation}
 D\ge\sigma\rho-\gamma-6mAB+6BR_++6AR_-.
 \label{ms:finite:countcore:physicalD}
\end{equation}
We now express this inequality in the common conditional cubic moments.

Now \(A=(p-a_0)/h\), \(B=(q-a_1)/h\). Substitution in
\eqref{ms:finite:countcore:physicalD}, followed by
\eqref{ms:finite:countcore:rhoidentity}--\eqref{ms:finite:countcore:gammaidentity}, yields
\begin{align}
 Dh^3\ge{}&B_\sigma(p)+6v f(a_0,a_1)\notag\\
 &+qw_0[3p(\sigma+3)-6a_0]
   +pw_1[3q(\sigma+1)-6a_1]\notag\\
 &-(\sigma+1)q\mu_0+(\sigma-1)p\mu_1,
 \label{ms:finite:countcore:Dhalfplane}\\
 f(a,b)&=a+b-a^2-b^2-ab,\qquad
 B_\sigma(p)=v[(\sigma-1)+2p-(2\sigma+6)v].\notag
\end{align}
This supplies another necessary halfplane in the two common cubic moments.
Inserting the one-sided cubic bounds from
\eqref{ms:finite:countcore:momentbounds} gives the stronger lower bound
\begin{align}
 Dh^3\ge{}&B_\sigma+6vf(a_0,a_1)\notag\\
 &+q\left\{[3p(\sigma+3)-(\sigma+7)a_0]w_0
                    +(\sigma+1)w_0^2/a_0\right\}\notag\\
 &+p\left\{[3q(\sigma+1)-(\sigma+5)a_1]w_1
                    +(\sigma-1)w_1^2/a_1\right\}.
 \label{ms:finite:countcore:Dsharp}
\end{align}
At \(a_i=0\), its conditional variance is zero and its contribution is
defined as zero. Lower substitution for a variance in a linear term
requires a nonnegative coefficient; otherwise a signed interval bound is
used. The concave polynomial \(f\) has its minimum on a clipped rectangle
at a vertex. On \(a+b\le d<1/2\) it is also increasing in each
coordinate. Any stronger monotone inversion of
\eqref{ms:finite:countcore:Dsharp} requires its own derivative-sign bound on the
whole interval; failure of that bound gives no exclusion.

\subsubsection{Strict outgoing bounds}

\begin{lemma}[A family of polynomial minorants]
\label{ms:finite:countcore:quartic}
For \(a,c>0\), put
\[
 Q_{a,c}(x)=x(ca-x)(x+a)^2,
 \qquad K_c=\max_{0\le t\le c}t(c-t)(t+1)^2.
\]
Then \(K_c>0\) and
\[
 Q_{a,c}(x)\le a^4K_c\mathbf1_{\{x>0\}}\qquad(x\in\mathbb R).
\]
If \(W\) is centered, of variance \(V\), third moment \(M_3\) and
fourth moment \(M_4\), then, writing
\(B_1=(c-2)a\), \(B_2=(2c-1)a^2\), \(B_3=ca^3\),
\begin{align}
 \Pp(W>u)\ge\frac1{a^4K_c}\bigl\{
 &-M_4+(B_1+4u)M_3+(B_2-3B_1u-6u^2)V\notag\\
 &-u^4-B_1u^3+B_2u^2-B_3u\bigr\}.
 \label{ms:finite:countcore:outgoingquartic}
\end{align}
\end{lemma}
\begin{proof}
The polynomial is nonpositive on \(x\le0\) and \(x\ge ca\).
On the remaining interval scale by \(a\). Expand \(Q_{a,c}(W-u)\)
and use \(\E W=0\).
\end{proof}

The positive maximizer defining \(K_c\) is
\[
 t_c=\frac{3c-2+\sqrt{9c^2+4c+4}}8.
\]
For \(c=2\), \(K_2=9/4+3\sqrt3/2\). There are also exact rational
normalizations. For \(c=3/2\),
\[
 2-t(3/2-t)(t+1)^2
 =(t-1)^2\bigl((t+5/4)^2+7/16\bigr).
\]
More generally, if \(t_*>0\),
\[
 c=\frac{2t_*(2t_*+1)}{3t_*+1},\qquad
 K=\frac{t_*^2(t_*+1)^3}{3t_*+1},
\]
then
\[
 K-t(c-t)(t+1)^2=(t-t_*)^2
 \left[\left(t+\frac{t_*^2+3t_*+1}{3t_*+1}\right)^2
       +\frac{t_*^2(2t_*^2+4t_*+1)}{(3t_*+1)^2}\right].
\]
These identities certify the pointwise inequality globally, rather than at
selected values of the random variable.

For \(W=W_k\), independence gives
\[
 M_3=(n-k)\mu_0+k\mu_1,\qquad
 M_4=(n-k)\E U_0^4+k\E U_1^4
       +3\{V_k^2-(n-k)w_0^2-kw_1^2\}.
\]
Using \eqref{ms:finite:countcore:momentbounds} makes the outgoing lower bound
affine in the \emph{same} pair \((\mu_0,\mu_1)\) used by the
incoming bounds. An alternative independent bound is
\[
 \E(U_i^4-4uU_i^3)
 \le w_i\max\{L_i^2+4uL_i,R_i^2-4uR_i\},
 \qquad U_i\in[-L_i,R_i].
\]
It follows by maximizing the convex quadratic \(x^2-4ux\) on the
support. A phase interval crossing an endpoint switch retains both cases.

For \(u<0\), Cantelli also gives
\(\Pp(W_k>u)\ge u^2/(V_k+u^2)\). A further elementary bound is
useful when a positive-class residual is forced upwards. For
\(T\in[-a,O]\), \(\E T=0\), \(\E T^2=w\), and \(-a\le b<O\),
\[
 \Pp(T>b)\ge\max\left\{0,
             \frac{w-ab}{(O+a)(O-b)}\right\}.
\]
Indeed \((T+a)(T-b)\) is nonpositive on \(T\le b\), and is at
most \((O+a)(O-b)\) otherwise. For the own count, choose a buffer
\(B\ge0\) and set \(b=u+B+(k-1)a_1\). Any one of the \(k\)
class-one variables exceeding \(b\) makes their sum exceed \(u+B\).
Independence gives the lower probability \(1-(1-r)^k\). The other
group is at least \(-B\) with probability at least
\[
 \max\left\{\frac{B^2}{(n-k)w_0+B^2},
             \frac{B}{B+(n-k)a_0}\right\}.
\]
The two groups are independent, so their lower probabilities may be
multiplied. At \(B=0\) use one only if the other-group variance is
zero, and otherwise zero. Thresholds outside \([-a,O)\) use the
corresponding certain or empty strict event. No atomic CDF is differentiated.

\subsubsection{Incoming exponential bounds}

For a centered \(Y\le B\), \(B>0\), of variance at most \(w\),
\begin{equation}
 \E e^{tY}\le1+\frac w{B^2}(e^{tB}-1-tB)
 \qquad(t\ge0).
 \label{ms:finite:countcore:bennett}
\end{equation}
The pointwise inequality follows because
\((e^{tx}-1-tx)/x^2=t^2\int_0^1(1-r)e^{trx}\,dr\) is increasing
in \(x\). Taking expectations proves \eqref{ms:finite:countcore:bennett}.

\begin{lemma}[Sharp one-sided MGF and its tangent]
\label{ms:finite:countcore:sharpmgf}
For the same variables,
\[
 \E e^{tY}\le M(w;B,t):=
 \frac{B^2e^{-tw/B}+we^{tB}}{B^2+w}.
\]
For fixed \(B,t\), the function \(M\) is increasing and concave in
\(w\ge0\). Hence every tangent \(A+D w\) bounds it above on that
whole half-line, with \(A\ge1\), \(D\ge0\).
\end{lemma}
\begin{proof}
For \(w>0\), put \(a=-w/B\). The quadratic matching \(e^{tx}\)
and its derivative at \(a\), and matching its value at \(B\),
majorizes \(e^{tx}\) for every \(x\le B\): its Hermite remainder
has sign \((x-a)^2(x-B)\le0\). Its quadratic coefficient is
nonnegative. The two-point law on \(\{a,B\}\) with masses
\(B^2/(B^2+w),w/(B^2+w)\) has mean zero and variance \(w\),
and therefore evaluates this majorant's expectation exactly. Variance
upper bounds remain valid. Differentiating gives
\[
 M'(w)=\frac{B^2}{(B^2+w)^2}
 \left[e^{tB}-e^{-tw/B}\left(1+\frac{t(B^2+w)}B\right)\right]\ge0.
\]
For \(x=t(B^2+w)/B\),
\[
 M''(w)=\frac{B^2e^{-tw/B}}{(B^2+w)^3}
        [x^2+2x+2-2e^x]\le0.
\]
At zero variance the law is concentrated at zero. Concavity proves the
tangent assertion, including its value at zero.
\end{proof}

Apply these bounds to \(-U_0\) and \(-U_1\), whose upper supports
are \(O_0,a_1\). Conditional independence bounds
\(\E e^{-tW_j}\) by the product of the two conditional MGF bounds,
with multiplicities \(n-j,j\). A zero upper support forces that
conditional variable to vanish and contributes the factor one.
The original common allocation may be retained. For example, tangents
to the sharp bounds give a product
\[
 (A_0+D_0c_0(1-\alpha))^{n-j}
 (A_1+D_1c_1\alpha)^j,
 \qquad c_0=\frac{\lambda_+}{n(1-p_+)},\quad
 c_1=\frac{\lambda_+}{np_-}.
\]
Its logarithm is concave. Its maximum on a closed allocation interval
is at an endpoint or its sole stationary point; zero-slope cases are
monotone or constant. The same \(\alpha\) is used in both factors.

If \(M\ge\E e^{tY}\), \(\E Y=0\), \(t,d>0\), then
\begin{align}
 \Pp(Y\ge d)&\le Me^{-td},\notag\\
 \Pp(Y\ge d)&\le
 \frac{e^\eta M-1-\eta}{e^{td+\eta}-1-td-\eta}
 \qquad(\eta\ge0).
 \label{ms:finite:countcore:shiftedexcess}
\end{align}
For the second inequality use the nonnegative function
\(e^{ty+\eta}-1-ty-\eta\), increasing on \(y\ge d\).
Its expectation is bounded by the numerator. The case \(\eta=0\)
is the centered exponential-excess bound. These are closed-tail bounds.
For \(j>k\), \(d=j-k-u>0\) on \(u\le1/2\); Cantelli provides
the independent bound \(V_j/(V_j+d^2)\).

\subsubsection{Incoming stop-loss comparisons}

\begin{lemma}[Two-point stop-loss comparison]
\label{ms:finite:countcore:stoploss}
Let \(\E Y=0\), \(\E Y^2\le w\), \(Y\le B\), with
\(w,B>0\). Let \(Z\) take \(-w/B,B\) with probabilities
\(B^2/(B^2+w),w/(B^2+w)\). For every real \(r\) and
\(a\in\{2,3\}\),
\[
 \E(Y-r)_+^a\le\E(Z-r)_+^a.
\]
The comparison law is the unit atom at zero if \(w=0\) or if
\(B=0\) and the centered original law is one-sided.
\end{lemma}
\begin{proof}
Put \(c=-w/B\). If \(r\ge B\), both positive parts vanish.
If \(c\le r<B\), the quadratic
\[
 P(x)=\frac{(B-r)^a}{(B-c)^2}(x-c)^2
\]
majorizes \((x-r)_+^a\) on \(x\le B\). For \(x>r\), this
follows by differentiating \((x-r)^a/(x-c)^2\); its derivative is
nonnegative for \(a=2,3\). Equality holds at \(c,B\).
If \(r<c\), use \(P(x)=(x-r)^2\) when \(a=2\). For \(a=3\),
put \(\delta=c-r>0\), \(\ell=B-c\), and use
\[
 P(x)=\delta^3+3\delta^2(x-c)+(3\delta+\ell)(x-c)^2.
\]
It is globally positive and
\(P(x)-(x-r)^3=(B-x)(x-c)^2\) on \(r\le x\le B\).
Every majorant has a nonnegative quadratic coefficient, so its
expectation increases when the variance is replaced by \(w\).
The comparison law has the specified first two moments and evaluates
the majorant exactly.
\end{proof}

Conditioning on the other summands and replacing them one at a time
extends the lemma to sums of independent centered variables. For the two
conditional types, the comparison sum is a convolution of two scaled
binomial laws. For every \(t>0\),
\begin{equation}
 \Pp(-W_j\ge d)
 \le\E[1+t(-W_j-d)]_+^a
 \le\E[1+t(Z_j-d)]_+^a,
 \qquad a=2,3.
 \label{ms:finite:countcore:stoplosstail}
\end{equation}
At \(d=j-k-u\), the last expression is convex and nondecreasing
in \(u\); its endpoint chord is an affine upper bound throughout a
closed phase interval. A fixed \(t\) is chosen for the whole chord.

For \(a=2\), its optimal value at one distance is obtained exactly.
For comparison atoms \(x_i\) of masses \(c_i>0\), write
\[
 F(t)=\sum_i c_i\max\{0,1+t(x_i-d)\}^2.
\]
This function is continuously differentiable and convex. Between its
positive breakpoints \(1/(d-x_i)\), the active sum is
\(C+2Bt+At^2\), \(A\ge0\). Check the breakpoints, the stationary
point \(-B/A\) when it lies in the active interval, and a possible
final constant interval. Since \(F'(0)=-2d<0\), the zero endpoint
is unnecessary. Repeated atoms and breakpoints are merged. Rational
input gives rational candidates and exact finite sums.

Common allocation can be preserved by a finite closed partition of
\([\alpha_-,\alpha_+]\). On a part \([a,b]\), use
\[
 w_0\le\frac{\lambda_+(1-a)}{n(1-p_+)},\qquad
 w_1\le\frac{\lambda_+b}{np_-},
\]
intersected with any prior valid variance caps. Form one matched phase
chord on each part. The line whose two endpoint values are the respective
maxima of those chords dominates every part on the entire phase interval.
It is therefore a whole-allocation bound. One may not join independently
optimized endpoint values unless they arise from such complete chords.

\subsubsection{Common cubic moments and phase coherence}

For \(t\ge0\), the function
\[
 r_t(x)=\frac{e^{tx}-1-tx}{x^2}
       =t^2\int_0^1(1-r)e^{trx}\,dr
\]
is nonnegative, increasing and convex, with its continuous value at zero.
On \([A,B]\), \(A<0<B\), its chord has coefficients
\[
 c_2=\frac{Br_t(A)-Ar_t(B)}{B-A}\ge0,\qquad
 c_3=\frac{r_t(B)-r_t(A)}{B-A}\ge0.
\]
Consequently
\[
 e^{tx}\le1+tx+c_2x^2+c_3x^3\quad(A\le x\le B).
\]
For \(Y=-U_i\), this gives an affine upper MGF in the original cubic
moment \(\mu_i\), with the third-moment sign reversed. Directed
coefficient enclosures give a globally nonnegative convex upper function
\[
 Q_i(\mu_i)=\max\{0,1+C_2w_{i,+}-C_{3,-}\mu_i,
                         1+C_2w_{i,+}-C_{3,+}\mu_i\},
\]
\[
 0\le C_{3,-}\le c_3\le C_{3,+},\quad c_2\le C_2.
\]
The maximum is essential when the sign of \(\mu_i\) is not fixed.

Choose fixed positive pivots \(P_0,P_1\). Weighted AM--GM gives the
convex common-moment upper bound
\begin{equation}
 Q_0^{n-j}Q_1^j
 \le G_j:=P_0^{n-j}P_1^j
 \left(\frac{n-j}{n}\frac{Q_0}{P_0}
                +\frac jn\frac{Q_1}{P_1}\right)^n.
 \label{ms:finite:countcore:convexmgf}
\end{equation}
The bracket is nonnegative and convex, and \(x^n\) is increasing and
convex on the nonnegative half-line. Thus \(G_j\) is convex in
\((\mu_0,\mu_1)\). At \(j=0,n\), omit zero powers.
Chernoff gives the incoming bound
\(e^{t_j(u-(j-k))}G_j\). Its exponential phase factor is bounded by
its positive endpoint chord. Alternatively use
\eqref{ms:finite:countcore:shiftedexcess} with the convex numerator
\(\max\{0,e^\eta G_j-1-\eta\}\). For \(x>0\),
\[
 \left(\frac1{e^x-1-x}\right)''
 =\frac{e^{2x}+(x-3)e^x+2}{(e^x-1-x)^3}>0.
\]
The numerator and its derivative vanish at zero, and its second derivative
is \(e^x(4e^x+x-1)>0\). Hence the reciprocal-excess phase factor
also has a positive upper chord. Each incoming strategy is fixed for the
whole moment polygon and whole phase interval.

Let \(\mathcal P\) be a bounded closed polygon containing the actual
pair \((\mu_0,\mu_1)\). The cubic moment bounds, the rho and signed
gamma strips from \eqref{ms:finite:countcore:rhoidentity}--\eqref{ms:finite:countcore:gammaidentity},
and \eqref{ms:finite:countcore:Dhalfplane} supply such necessary halfplanes.
In a probability box of midpoint \(p_*\) and halfwidth \(\delta\),
replacing \(-q\mu_0+p\mu_1\) by its midpoint coefficients costs at
most \(\delta\max|\mu_0+\mu_1|\); the signed gamma strip costs at
most \(\delta\max|\mu_1-\mu_0|\). All remaining factors receive
outward interval bounds. These enlargements preserve the actual pair.
An empty polygon is a genuine necessary-condition contradiction.

For a polynomial of degree at most \(d\) on \([l,r]\), write
\[
 P(u)=\sum_{i=0}^d P_i\binom di t^i(1-t)^{d-i},
 \qquad t=(u-l)/(r-l).
\]
The Bernstein basis is nonnegative and sums to one, so
\(P(u)\le\max_iP_i\). An affine phase chord has coefficients its
values at \(l+(r-l)i/d\). Combining these positive coefficients
with \eqref{ms:finite:countcore:convexmgf} and the affine common-moment outgoing
quartic produces convex coefficient functions \(F_i(\mu_0,\mu_1)\).
The degree-four own polynomial needs \(d=4\). Constant intervals are
evaluated directly.

\begin{lemma}[Vertex reduction with normalization]
\label{ms:finite:countcore:vertex}
If \(F\) is convex on a compact polygon and \(L\) is affine and
strictly positive there, then
\[
 \frac{F(x)}{L(x)}\le\max_{v\ \mathrm{vertex}}\frac{F(v)}{L(v)}.
\]
If the acquired rho floor is
\(\max\{r_0,L(\mu_0,\mu_1)\}\), \(r_0>0\), split the polygon
along \(L=r_0\) and use the vertices of both closed pieces.
\end{lemma}
\begin{proof}
Write \(x=\sum_v\theta_vv\). Convexity gives
\(F(x)\le\sum_v\theta_vF(v)\), which is at most
\(\max_v(F(v)/L(v))\sum_v\theta_vL(v)\). The latter sum is
\(L(x)\). The same proof applies to segments and single points.
\end{proof}

One must maximize over all controls, normalization pieces and their
vertices before minimizing over a finite menu of rates, pivots or quartic
shapes. A pointwise minimum of convex functions need not be convex.
For the same reason, replacing a convex incoming bound pointwise by its
minimum with one is not part of this vertex argument. A constant bound
may instead be chosen as a whole-polygon strategy.

\subsubsection{From a count bound to the normalized score}

\begin{theorem}[Whole-box normalized count bound]
\label{ms:finite:countcore:scoretheorem}
Fix a closed box of acquired parameters and a phase interval \([l,r]\).
Let \(B_k\) bound the binomial CDF above, \(b_k^-\) the own mass
below, and \(b_j^+\) the incoming masses above, throughout that box.
Choose one outgoing lower bound \(L_k(u,\mu)\) proved above and
incoming nonnegative upper bounds \(A_j(u,\mu)\) proved above.
Let \(N(u)\) be a lower bound for
\[
 \Phi\left(\frac{k+u-np}{\sqrt n\sqrt{v+\lambda/n}}\right)
\]
throughout the same box. Suppose their combination has certified convex
Bernstein upper controls \(F_i(\mu)\):
\[
 B_k-b_k^-L_k(u,\mu)+\sum_{j>k}b_j^+A_j(u,\mu)-N(u)
 \le\max_i F_i(\mu).
\]
This control is obtained by the polynomial, affine and convex constructions
above, not assumed as a probability inequality. On each normalization
piece let \(r(\mu)>0\) be its affine lower rho bound. If \(n_+\ge\sqrt n\),
the positive normalized CDF error on the whole box is at most
\begin{equation}
 n_+\max_{i,\,v\ \mathrm{vertex}}
                  \frac{\max\{0,F_i(v)\}}{r(v)}.
 \label{ms:finite:countcore:scorebound}
\end{equation}
The maximum includes every closed normalization piece.
\end{theorem}
\begin{proof}
Insert the outgoing and incoming probability bounds into
\eqref{ms:finite:countcore:countineq}. For the outgoing term,
\(b_k\Pp(W_k>u)\ge b_k^-\Pp(W_k>u)\ge b_k^-L_k\),
even if the polynomial lower bound is negative. Incoming upper bounds
are nonnegative, so the upper mass substitutions are valid. Subtract
the normal lower bound, take the positive part, and use \(\rho\ge r\).
The positive part of a convex function is convex. Apply
Lemma~\ref{ms:finite:countcore:vertex}, and only then enlarge \(\sqrt n\).
\end{proof}

A convenient normal lower envelope uses a lower value \(N_l\) at the
left phase endpoint and a uniform lower derivative \(g\ge0\).
For \(h_-\le\sqrt{v+\lambda/n}\le h_+\), put
\[
 X_*=\frac{\max|k+u-np|}{\sqrt n\,h_-},\qquad
 g\le\frac{e^{-X_*^2/2}}{\sqrt{2\pi n}\,h_+}.
\]
Then \(N(u)=N_l+g(u-l)\) is valid. Only the smooth Gaussian is
differentiated. Cantelli's smooth incoming envelope has derivative
\(2Vd/(V+d^2)^2\), whose maximum over a distance interval occurs at
an endpoint or \(d=\sqrt{V/3}\). The exponential, centered and shifted
envelopes likewise have explicit upper derivatives or convex endpoint
chords. A matched pair of left value and slope is always retained.
One affine envelope dominates another on an interval precisely when it
does so at both endpoints. These observations justify retaining phase
cancellation rather than maximizing each term at a different phase.

\subsubsection{Complete roots and finite certificates}

If acquired \(p,h,H\) bounds imply
\[
 np_- -h_+H_-\ge-\tfrac12,\qquad
 np_+ +h_+H_+\le n+\tfrac12,
\]
then every allowed threshold has a nearest count \(k=0,\ldots,n\)
and \(u\in[-1/2,1/2]\). A tie belongs to both closed roots.
The original product roots are
\[
 [p_-,p_+]\times[0,\Lambda]\times[0,1]
                  \times[-1/2,1/2]\times[0,d_0],
 \qquad k=0,\ldots,n.
\]
Necessary variance, moment, support and inward contractions act only
inside these roots. A proof must still cover the original products.
If thresholds outside this range are possible, their tails need separate
bounds; they are not discarded by the nearest-count notation.

\begin{proposition}[Finite closed-cover implication]
\label{ms:finite:countcore:forest}
Suppose every original root has a finite binary partition by rational
strict interior cuts. Retain both closed children at each cut. If every
leaf either has a proved necessary-condition contradiction or has the
whole-box upper bound \eqref{ms:finite:countcore:scorebound} below a fixed
\(c<\CE\), then every allowed positive normalized CDF error is below
\(\CE\). This conclusion uses all \(n+1\) original roots.
\end{proposition}
\begin{proof}
A strict cut partitions a closed box into two closed boxes whose union
is the parent. Finite induction up each tree proves its root statement,
including all seams. Union over the count roots proves the assertion.
Zero-variance faces use deterministic conditional laws as above; equality
in any necessary-condition interval test is not an emptiness certificate.
\end{proof}

Rational arithmetic implements this proposition over continuous domains.
For a square root, integer squaring brackets the positive radical;
polynomial identities and Bernstein coefficients are checked over exact
rationals. Exponentials and the Gaussian CDF are enclosed by finite series
with signed or absolute remainder bounds. Every coefficient is rounded in
the direction used in its proved inequality. The mathematical content is
the whole-box inequality and the finite cover, not a sample of parameter
values or a floating-point optimum.

Finally, this section bounds the positive closed CDF branch. For a negative
Kolmogorov discrepancy use the whole reflected marginal and
\[
 \Phi(x)-\Pp(S_n<x)=\Pp(-S_n\le-x)-\Phi(-x).
\]
Resources for that oriented law must be acquired in the same manner.
Treating both signs of the threshold inside one row does not by itself
reflect the marginal or its signed moments.



\subsection{One summand: a closed cubic-dual implication}
\label{ms:finite:n1}
\begin{condition}[One-summand dual cover]\label{ms:finite:n1-cover}
A finite closed cover of $[-2,1]$ supplies, for each rational cell $[a,b]$,
rationals $u,v,L_a$ such that $L_a\le\Phi(a)$ and
\[
 L_a-v+ux+vx^2+\tfrac25|x|^3-\ind_{x\le b}\ge0
 \qquad(x\in\R).
\]
The polynomial inequalities are established on the whole real line,
including both infinite rays and the closed event boundary. The prescribed
cover uses 150 cells.
\end{condition}
\begin{proposition}\label{ms:finite:n1-bound}
Under Condition~\ref{ms:finite:n1-cover}, every standardized law with
finite third absolute moment satisfies $\Delta_1(P)\le(2/5)\rho(P)$.
\end{proposition}
\begin{proof}
Taking expectations cancels $u\E X$ and $-v+v\E X^2$, giving
$F_X(b)-L_a\le(2/5)\rho$. For $t\in[a,b]$, monotonicity of both
CDFs gives $F_X(t)-\Phi(t)\le F_X(b)-L_a$.
For $t\le-2$, Cantelli gives $F_X(t)\le1/(1+t^2)\le1/5$.
For $t\ge1$, $1-\Phi(t)\le1-\Phi(1)<2/5$; for example,
$\phi(0)>3/8$ and $e^{-s^2/2}\ge1-s^2/2$ on $[0,1]$ imply
$\Phi(1)>1/2+(3/8)(5/6)=13/16$. Since $\rho\ge1$, these tails
suffice. For the opposite discrepancy,
\[
 \Phi(t)-F_X(t)\le\Phi(t)-F_X(t-)
   =F_{-X}(-t)-\Phi(-t).
\]
Apply the same result to the reflected law, retaining all atoms.
\end{proof}

\subsection{Two summands: a rational tensor certificate}\label{ms:finite:n2:section}
\subsubsection{Result and proof boundary}
For a standardized law $P$, write
\[
 \E X=0,\qquad \E X^2=1,\qquad \rho(P)=\E|X|^3<\infty,
\]
\[
 \Delta_n(P)=\sup_x\left|P^{*n}(( -\infty,x\sqrt n])-\Phi(x)\right|,
 \qquad C_n(P)=\frac{\sqrt n\,\Delta_n(P)}{\rho(P)},\qquad C_n=\sup_P C_n(P).
\]
Let $\phi=\Phi'$ and $\phi_0=(2\pi)^{-1/2}$. Set
\[
 K=\frac{2\sqrt2}{7}=0.404061017820884\ldots,
 \qquad \CE=\frac{3+\sqrt{10}}{6\sqrt{2\pi}}=0.409732183702396\ldots.
\]
\begin{theorem}[Certificate-conditional two-summand bound]\label{ms:finite:n2:thm:main}
If the finite positivity and coverage obligations specified below are
verified, then for every standardized law $P$ with finite third absolute moment,
\[
 \boxed{\Delta_2(P)\le\frac27\rho(P).}
\]
Equivalently, $C_2\le K<\CE$.
\end{theorem}
The analytic reduction below gives the stated upper bound conditional on
exact verification of the finite positivity assertions.

The sole external numerical probability estimate is Shevtsova's published
structural inequality \cite{Shevtsova2013}:
\begin{equation}\label{ms:finite:n2:eq:struct}
 \frac{\sqrt n\,\Delta_n(P)}{\rho(P)}
 \le a\left(1+\frac b{\rho(P)}\right),
 \qquad a=\frac{1661}{5000},\quad b=\frac{429}{1000}.
\end{equation}
It implies
\begin{equation}\label{ms:finite:n2:eq:large}
 \rho\ge2\quad\Longrightarrow\quad
 C_n(P)\le a(1+b/2)<K.
\end{equation}
The last comparison is rational after squaring, since $K^2=8/49$.

\subsubsection{Reduction of a putative maximizer to \texorpdfstring{$[-6,6]$}{[-6,6]}}
\begin{lemma}[Attainment under the contrary hypothesis]\label{ms:finite:n2:lem:attain}
If $C_2>K$, the supremum is attained by a standardized law $P$ with $\rho(P)<2$.
\end{lemma}
\begin{proof}
Lemma~\ref{ms:attainment}, applied at $n=2$, gives an unrestricted
maximizing law and a finite threshold with positive closed-CDF discrepancy. If its third
absolute moment were at least $2$, \eqref{ms:finite:n2:eq:large}
would give $C_2<K$, contrary to the hypothesis.
\end{proof}

For the positive closed-CDF maximizing pair supplied by that lemma, write
\[
 \begin{gathered}
 t=x\sqrt2,\qquad F=P^{*2}(( -\infty,t]),\qquad D=F-\Phi(x)>0,\\
 C=\frac{\sqrt2D}{\rho}=C_2,\qquad
 \lambda=\frac D\rho,\qquad m=\E[X|X|].
 \end{gathered}
\]
Here $m$ denotes the signed second moment $\beta$ of
Lemma~\ref{ms:contact}, and $|m|\le1$.

\begin{lemma}[Explicit contact inequality]\label{ms:finite:n2:lem:contact}
Define
\begin{align}
 Q(y)={}&\lambda|y|^3-\frac{x\phi(x)+3D}{2}y^2
       -(\sqrt2\phi(x)+3\lambda m)y\notag\\
       &\hspace{20mm}+2F+\frac{x\phi(x)+D}{2}.\label{ms:finite:n2:eq:Q}
\end{align}
Then
\begin{equation}\label{ms:finite:n2:eq:contact}
 2P(( -\infty,t-y])\le Q(y)\quad(y\in\R),
 \qquad Q(y)=2P(( -\infty,t-y])\quad(y\in\supp P).
\end{equation}
\end{lemma}
\begin{proof}
Apply Lemma~\ref{ms:contact} at $n=2$, $H=t$ and $z=x$.
In \eqref{ms:contact-majorant}, substitute $\beta=m$,
$C/\sqrt2=\lambda$ and $C\rho/\sqrt2=D$. The polynomial displayed
here is $2Q_t$, so the asserted inequality and equality on the entire
closed support follow, including all atoms.
\end{proof}

\begin{lemma}[Six-unit support bound]\label{ms:finite:n2:lem:support}
Under $C_2>K$, every maximizing witness in Lemma~\ref{ms:finite:n2:lem:attain} satisfies
$\supp P\subset(-6,6)$.
\end{lemma}
\begin{proof}
Equation \eqref{ms:finite:n2:eq:struct} gives $C<1/2$ and $C\rho<a(2+b)<1$.
Hence $\lambda>2/7$, $D<1/\sqrt2<5/7$, and $\lambda<5/14$.
Using $\phi(x)<2/5$, $|x\phi(x)|<1/4$, $\sqrt2<10/7$,
$|m|\le1$, and $F,D\ge0$, we obtain
\[
 Q(y)\ge P_6(|y|),\qquad
 P_6(r)=\frac27r^3-\frac{67}{56}r^2-\frac{23}{14}r-\frac18.
\]
The exact identity
\[
 P_6(6+u)-2
 =\frac27u^3+\frac{221}{56}u^2+\frac{104}{7}u+\frac{373}{56}>0
 \quad(u\ge0)
\]
contradicts $Q\le2$ on the support whenever $|y|\ge6$.
\end{proof}
No finite-support assumption appears in this reduction.

\subsubsection{The two moment parameters}
Let $\gamma=\E X^3$. Every standardized law satisfies
\begin{equation}\label{ms:finite:n2:eq:momentregion}
 1\le\rho,\qquad |\gamma|\le\rho,\qquad
 \rho^2(\gamma^2+4)-(\gamma^2+2)^2\ge0.
\end{equation}
The last inequality is Assumption~\ref{ms:finite:inputs}, item~\ref{ms:finite:input:orig:thm:defect}; squaring its positive
sides gives \eqref{ms:finite:n2:eq:momentregion}. In the remaining problem $\rho\le2$,
so it suffices to consider the compact domain
\begin{equation}\label{ms:finite:n2:eq:D}
 \mathcal D=\{(\gamma,\rho)\in[-2,2]\times[1,2]:
 \rho^2(\gamma^2+4)-(\gamma^2+2)^2\ge0\}.
\end{equation}
These necessary moment constraints give an outer domain for the certificate.

\subsubsection{The tensor majorants}
Define
\[
 v_0(y)=1,\quad v_1(y)=y,\quad v_2(y)=y^2,
 \quad v_3(y)=y^3,\quad v_4(y)=|y|^3,
 \qquad \mu=(1,0,1,\gamma,\rho).
\]
For rational coefficients $c=(c_{ij})_{0\le i\le j\le4}$ set
\begin{align}
 M_c(u,v)&=\sum_i c_{ii}v_i(u)v_i(v)
 +\sum_{i<j}c_{ij}\{v_i(u)v_j(v)+v_j(u)v_i(v)\},\label{ms:finite:n2:eq:tensor}\\
 q_c(\gamma,\rho)&=\sum_i c_{ii}\mu_i^2
                 +2\sum_{i<j}c_{ij}\mu_i\mu_j.\label{ms:finite:n2:eq:expect}
\end{align}
Independence and the common distribution give the exact identity
\begin{equation}\label{ms:finite:n2:eq:factor}
                      \E M_c(X_1,X_2)=q_c(\gamma,\rho).
\end{equation}
On each quadrant $M_c$ is a polynomial of degree at most three in each
coordinate. Its expectation is a polynomial of total degree at most two
in $(\gamma,\rho)$.

For $j=0,\ldots,69$, put
\[
 a_j=-\frac52+\frac j{20},\qquad b_j=a_j+\frac1{20}.
\]
The certificate condition requires one rational vector $c^{(j)}$ for each interval and the inequalities
\begin{align}
 M_{c^{(j)}}(u,v)&\ge\mathbf1_{\{u+v\le b_j\}}
                 &&(u,v\in[-6,6]),\label{ms:finite:n2:eq:C1}\\
 q_{c^{(j)}}(\gamma,\rho)&\le\Phi(a_j/\sqrt2)+\frac27\rho
                 &&((\gamma,\rho)\in\mathcal D).\label{ms:finite:n2:eq:C2}
\end{align}
This prescribed certificate format has 70 threshold rows and 1050 coefficients;
coefficients are ordered lexicographically by $(i,j)$ with $i\le j$.

\begin{lemma}[Certificate implication]\label{ms:finite:n2:lem:certimp}
Equations \eqref{ms:finite:n2:eq:C1}--\eqref{ms:finite:n2:eq:C2} imply
$\Delta_2(P)\le2\rho(P)/7$ for every standardized law on $[-6,6]$ with
$\rho\le2$.
\end{lemma}
\begin{proof}
If $t\in[a_j,b_j]$, monotonicity, \eqref{ms:finite:n2:eq:factor}, and the certificate give
\begin{align*}
 P(X_1+X_2\le t)-\Phi(t/\sqrt2)
 &\le\E M_{c^{(j)}}(X_1,X_2)-\Phi(a_j/\sqrt2)\\
 &\le\frac27\rho.
\end{align*}
For $t\le-5/2$, Cantelli's inequality gives
$P(X_1+X_2\le t)\le2/(2+t^2)\le8/33<2/7\le2\rho/7$.
For $t\ge1$, $1-\Phi(t/\sqrt2)\le1-\Phi(1/\sqrt2)<2/7$.
The last strict inequality is checked by the rational Gaussian enclosure below.
Thus the positive branch is bounded at every real threshold.
Applying the same result to $-X$ bounds
$\Phi(t/\sqrt2)-P(X_1+X_2<t)$, and hence also the negative right-cdf error.
Atoms and exact sum collisions are included, not removed by genericity.
\end{proof}

\subsubsection{Exact checking of the continuous domains}
\paragraph{Piecewise-polynomial positivity}
Each of the four quadrants of $[-6,6]^2$ is intersected with the two closed
half-planes $u+v\le b_j$ and $u+v\ge b_j$. Each resulting rational polygon is
triangulated. On each triangle, it suffices to prove nonnegativity of
$M_{c^{(j)}}-1$ on the first side or $M_{c^{(j)}}$ on the second.
The shared boundary belongs to both closed regions.

Write a degree-six polynomial in barycentric Bernstein form on a triangle:
\[
 p=\sum_{\alpha_0+\alpha_1+\alpha_2=6}
 B_\alpha\frac{6!}{\alpha_0!\alpha_1!\alpha_2!}
 \lambda_0^{\alpha_0}\lambda_1^{\alpha_1}\lambda_2^{\alpha_2}.
\]
If every $B_\alpha\ge0$, then $p\ge0$ on the whole triangle. Otherwise the
longest edge is bisected and both children are checked. A resource limit or
negative vertex value is a rejection, never an acceptance.

One exact procedure constructs the coefficients from a monomial expansion
and subdivides them by rational de Casteljau identities. A second procedure
reconstructs the homogeneous polynomial directly on \emph{every} child.
If $D_c$ clears polynomial denominators and $D_v$ clears vertex denominators,
write $U=D_vu$, $V=D_vv$, and $S=\lambda_0+\lambda_1+\lambda_2$.
For $p(u,v)=\sum a_{rs}u^rv^s$, the second construction expands
\[
 \sum_{r,s}(D_ca_{rs})D_v^{6-r-s}U^rV^sS^{6-r-s}.
\]
All its coefficients are integers. They are nonnegative exactly when the
corresponding Bernstein coefficients are nonnegative, because the omitted
denominators and multinomial factors are positive. This gives a direct
coefficient calculation without using the subdivision recurrence.

\paragraph{Moment-parameter positivity}
The scalar expression
$\Phi(a_j/\sqrt2)+2\rho/7-q_{c^{(j)}}(\gamma,\rho)$ is enclosed from below
on rational rectangles using degree-two Bernstein coefficients in each
variable. A rectangle outside $\mathcal D$ is discarded. This test is exact:
if $g_*$ is the minimum of $\gamma^2$ on its $\gamma$ interval and $r_+$ is
its upper $\rho$ endpoint, the rectangle is outside whenever
\[
 r_+^2(g_*+4)-(g_*+2)^2<0.
\]
Indeed, $r^2(w+4)-(w+2)^2$ increases in $r\ge0$ and decreases in
$w\ge0$ for $r\le2$. All other rectangles are certified or subdivided.

\paragraph{The normal cdf}
For $|t|\le5/2$,
\begin{equation}\label{ms:finite:n2:eq:normal}
 \Phi(t/\sqrt2)=\frac12+\frac1{2\sqrt\pi}
 \sum_{k=0}^\infty\frac{(-1)^k t^{2k+1}}{4^k k!(2k+1)}.
\end{equation}
For $t\ge0$ the term magnitudes decrease: their successive ratio is at most
$(25/16)(2k+1)/((k+1)(2k+3))<1$. Alternating partial sums therefore give
rational enclosures after bounding $1/(2\sqrt\pi)$. Machin's arctangent
identity and integer square roots supply those bounds. The second procedure
uses different rational prefactor bounds and different truncation orders.
Negative arguments are handled by symmetry. No platform implementation of
$\Phi$, exponential, sine, or cosine is used for acceptance.

\paragraph{The complete finite certificate}
The supplied certificate has exactly seventy consecutive threshold
rows and \(1050\) rational coefficients. Complete exact verification
establishes both the pointwise majorant and moment-parameter inequalities
on their stated closed domains. Either of the two positivity procedures
above can establish these finite obligations.

\subsubsection{Completion of the two-summand proof}
\begin{proof}[Proof of Theorem~\ref{ms:finite:n2:thm:main}]
Suppose $C_2>K$. Lemma~\ref{ms:finite:n2:lem:attain} supplies a global maximizing law
with $\rho<2$. Lemma~\ref{ms:finite:n2:lem:support} confines it to $(-6,6)$.
The assumed complete rational verification establishes \eqref{ms:finite:n2:eq:C1}--\eqref{ms:finite:n2:eq:C2};
Lemma~\ref{ms:finite:n2:lem:certimp} then gives $C_2(P)\le K$, a contradiction.
Thus $C_2\le K$. Finally, exact rational enclosures give $K<0.405<0.409<\CE$
using rational bounds and squaring.
\end{proof}
The bounded-support theorem in Lemma~\ref{ms:finite:n2:lem:certimp} treats arbitrary
probability measures, not merely finitely supported measures or an assumed
number of atoms. The passage to unrestricted laws uses global extremality,
not a truncation approximation.


\subsection{Three summands: independence-null certificates}\label{ms:finite:n3:section}
\subsubsection{Result and external input}
For a standardized law $P$, let
\[
 \E X=0,\quad \E X^2=1,\quad \rho(P)=\E|X|^3<\infty,
\]
\[
 \Delta_n(P)=\sup_x\left|P^{*n}(( -\infty,x\sqrt n])-\Phi(x)\right|,
 \qquad C_n=\sup_P\frac{\sqrt n\,\Delta_n(P)}{\rho(P)}.
\]
Write $\phi=\Phi'$ and set
\[
 \Lambda=59/250,\qquad K=\sqrt3\Lambda
 =0.408763990586255\ldots,
 \qquad \CE=\frac{3+\sqrt{10}}{6\sqrt{2\pi}}.
\]
\begin{theorem}[Certificate-conditional three-summand bound]\label{ms:finite:n3:thm:main}
If the finite positivity and coverage obligations specified below are
verified, then every standardized law with finite third absolute moment satisfies
\[
 \boxed{\Delta_3(P)\le\Lambda\rho(P).}
\]
Consequently $C_3\le59\sqrt3/250<\CE$.
\end{theorem}
The analytic reduction below gives the stated upper bound conditional on
exact verification of the polynomial inequalities over their complete domains.

The sole external numerical probability bound is Shevtsova's structural
estimate \cite{Shevtsova2013}:
\begin{equation}\label{ms:finite:n3:eq:struct}
 C_n(P)\le A\left(1+\frac B{\rho(P)}\right),
 \qquad A=\frac{1661}{5000},\quad B=\frac{429}{1000}.
\end{equation}
It implies
\begin{equation}\label{ms:finite:n3:eq:large}
 \rho(P)\ge(15/8)\quad\Longrightarrow\quad
 C_n(P)\le\frac{159456}{390625}<K.
\end{equation}
The last inequality is rational after squaring. No result about the optimal
two-point constant, no near-lattice transfer, and no effective eventual
threshold is used.

\subsubsection{A bounded witness would attain any violation}
\begin{lemma}[Attainment under the contrary hypothesis]\label{ms:finite:n3:lem:attain}
If $C_3>K$, a standardized law $P$ with $\rho(P)<15/8$ attains $C_3$.
\end{lemma}
\begin{proof}
Lemma~\ref{ms:attainment}, applied at $n=3$, gives an unrestricted
maximizing law and a finite threshold with positive closed-CDF discrepancy. If
$\rho(P)\ge15/8$, \eqref{ms:finite:n3:eq:large} would give
$C_3<K$, contrary to the hypothesis.
\end{proof}
For the positive closed-CDF pair supplied by that lemma, put
\[
 \begin{gathered}
 t=x\sqrt3,\quad F=P^{*3}(( -\infty,t]),\quad D=F-\Phi(x)>0,\\
 C=C_3=\frac{\sqrt3D}{\rho},\quad \lambda=\frac D\rho,
 \quad m=\E[X|X|],\quad |m|\le1.
 \end{gathered}
\]

\begin{lemma}[Explicit contamination contact inequality]\label{ms:finite:n3:lem:contact}
Define
\begin{align}
 Q(y)={}&\lambda|y|^3-\frac{x\phi(x)+3D}{2}y^2
 -(\sqrt3\phi(x)+3\lambda m)y \notag\\
 &\hspace{22mm}+3F+\frac{x\phi(x)+D}{2}.
 \label{ms:finite:n3:eq:Q}
\end{align}
Then
\begin{equation}\label{ms:finite:n3:eq:contact}
 3P^{*2}(( -\infty,t-y])\le Q(y)\qquad(y\in\R),
\end{equation}
with equality throughout $\supp P$.
\end{lemma}
\begin{proof}
Apply Lemma~\ref{ms:contact} at $n=3$, $H=t$ and $z=x$.
In \eqref{ms:contact-majorant}, substitute $\beta=m$,
$C/\sqrt3=\lambda$ and $C\rho/\sqrt3=D$. The polynomial displayed
here is $3Q_t$. The closed-CDF inequality and equality throughout the
closed support are therefore the corresponding conclusions of that lemma.
\end{proof}

\begin{lemma}[Uniform support bound]\label{ms:finite:n3:lem:support}
A violating maximizer in Lemma~\ref{ms:finite:n3:lem:attain} satisfies
$\supp P\subset(-6,6)$.
\end{lemma}
\begin{proof}
The structural bound and $1\le\rho<15/8$ give
\[
 \frac25<C<\frac{19}{40},\qquad C\rho<\frac{77}{100}.
\]
Using $12/7<\sqrt3<7/4$ and $\sqrt3>17/10$, we obtain
\[
 D<\frac9{20},\qquad \frac8{35}<\lambda<\frac27.
\]
Also $\phi(x)<2/5$ and $|x\phi(x)|<1/4$. The first follows from
$\pi>25/8$, and the second from the maximum
$(2\pi e)^{-1/2}<1/4$. Thus \eqref{ms:finite:n3:eq:Q} implies, with $r=|y|$,
\[
 Q(y)\ge\frac8{35}r^3-\frac45r^2-\frac85r-\frac18.
\]
For $r=6+u$, the difference of the right side and $3$ equals
\[
 \frac8{35}u^3+\frac{116}{35}u^2+\frac{472}{35}u+
 \frac{2197}{280}>0\qquad(u\ge0).
\]
But \eqref{ms:finite:n3:eq:contact} is an equality on the support, and its left side is
at most $3$. This excludes every support point with $|y|\ge6$.
\end{proof}
\begin{remark}
This is an extremizer reduction, not a truncation of arbitrary laws.
The original theorem assumes only a finite third absolute moment.
Higher moments are used below only after the hypothetical violating
maximizer has been proved bounded.
\end{remark}

\subsubsection{The admissible two-dimensional moment region}
Write $\gamma=\E X^3$. Assumption~\ref{ms:finite:inputs}, item~\ref{ms:finite:input:orig:thm:defect} gives
\[
 \rho\ge(\gamma^2+2)/\sqrt{\gamma^2+4}.
\]
Together with $|\gamma|\le\rho$ and $\rho\ge1$, every relevant law belongs to
\begin{equation}\label{ms:finite:n3:eq:momdomain}
 \mathcal D=\left\{(\gamma,\rho):
 |\gamma|\le\frac{15}{8},\quad1\le\rho\le\frac{15}{8},\quad
 \rho^2(\gamma^2+4)-(\gamma^2+2)^2\ge0\right\}.
\end{equation}
This is an outer domain: every point need not be realizable by a law.
The certificate checks all of it, so non-realizable points cannot cause
an invalid acceptance.

\subsubsection{Independence-null controls}
For a single coordinate define the features
\begin{equation}\label{ms:finite:n3:eq:features}
 (v_0,v_1,v_2,v_3,v_4,v_5,v_6)(y)
 =(1,y,y^2-1,y^3,|y|^3,y^4,y^5).
\end{equation}
Let $\mathcal I$ consist of the nondecreasing triples
$\alpha=(i,j,k)$ with $0\le i\le j\le k\le6$ such that either
$k<5$, or at least one entry is $1$ or $2$. There are $59$ such triples,
ordered lexicographically in the certificate files.
For each $\alpha$, let
\[
 V_\alpha(x_1,x_2,x_3)
 =\frac1{|\operatorname{Perm}(\alpha)|}
   \sum_{(i,j,k)\in\operatorname{Perm}(\alpha)}
        v_i(x_1)v_j(x_2)v_k(x_3),
\]
where the sum is over the \emph{distinct} permutations of the index triple.
For $c\in\mathbb Q^{59}$ put
\begin{equation}\label{ms:finite:n3:eq:majorant}
 M_c=\sum_{\alpha\in\mathcal I}c_\alpha V_\alpha.
\end{equation}

\begin{lemma}[Exact expectation without higher-moment parameters]\label{ms:finite:n3:lem:expect}
For independent standardized identically distributed bounded variables,
\begin{equation}\label{ms:finite:n3:eq:expect}
 \E M_c(X_1,X_2,X_3)=q_c(\gamma,\rho)
 :=\sum_{\substack{\alpha\in\mathcal I\\
                    \text{entries only }0,3,4}}
        c_\alpha\gamma^{\#3(\alpha)}\rho^{\#4(\alpha)}.
\end{equation}
In particular, $q_c$ has total degree at most three and has no unknown
fourth- or fifth-moment parameters.
\end{lemma}
\begin{proof}
The expectations of $v_1$ and $v_2$ are zero. Independence makes every
product containing either factor have expectation zero. By the definition
of $\mathcal I$, every product involving $v_5$ or $v_6$ contains such a
zero-mean factor in another coordinate. The only remaining products use
$v_0,v_3,v_4$, whose means are $1,\gamma,\rho$.
\end{proof}
This adds $24$ coefficients to the $35$ low-feature tensor coefficients.
In the resulting family, $49$ basis elements have zero expectation and
only $10$ contribute monomials in $(\gamma,\rho)$.
For example, all permutations of
\[
 X_1^4X_2,
 \qquad X_1^5(X_2^2-1),
 \qquad X_1^5X_2|X_3|^3
\]
have mean zero. These are exact independence identities, not assumptions
that the fourth or fifth moment has a prescribed value.

The general principle is useful beyond three summands. For any integrable
function of the other coordinates,
\[
 \E\{X_i A(X_1,\ldots,\widehat X_i,\ldots,X_n)\}=0,
\]
\[
 \E\{(X_i^2-1)B(X_1,\ldots,\widehat X_i,\ldots,X_n)\}=0.
\]
Such terms may reshape a pointwise majorant without changing its i.i.d.
expectation. For the present proof all functions are polynomials on a
bounded cube, so integrability is automatic.

\subsubsection{Finite-certificate implication}
\begin{condition}[Continuous-domain three-summand certificate]\label{ms:finite:n3:prop:certificate}
For each row $j$ of the supplied rational certificate, let $[a_j,b_j]$
be its threshold interval and $c^{(j)}\in\mathbb Q^{59}$ its coefficients.
The rows form an exact consecutive partition of $[-7/2,3/2]$ and satisfy
\begin{align}
 M_{c^{(j)}}(u,v,w)&\ge\ind_{\{u+v+w\le b_j\}}
       &&((u,v,w)\in[-7,7]^3),\label{ms:finite:n3:eq:pointwise}\\
 q_{c^{(j)}}(\gamma,\rho)&\le
      \Phi(a_j/\sqrt3)+\Lambda\rho
       &&((\gamma,\rho)\in\mathcal D).\label{ms:finite:n3:eq:scalar}
\end{align}
Section~\ref{ms:finite:n3:sec:check} gives sufficient exact tests for these
continuous-domain inequalities.
\end{condition}

\begin{proof}[Proof of Theorem~\ref{ms:finite:n3:thm:main}]
Suppose $C_3>K$. Lemmas~\ref{ms:finite:n3:lem:attain}--\ref{ms:finite:n3:lem:support} give a globally
maximizing violating law on $(-6,6)$, with $\rho<15/8$, and a maximizing
raw threshold $t$ with positive closed-CDF discrepancy.
If $t\in[a_j,b_j]$, then independence and Condition~\ref{ms:finite:n3:prop:certificate}
give
\begin{align*}
 P(X_1+X_2+X_3\le t)-\Phi(t/\sqrt3)
 &\le\E M_{c^{(j)}}(X_1,X_2,X_3)-\Phi(a_j/\sqrt3)\\
 &\le\Lambda\rho.
\end{align*}
The larger certificate cube $[-7,7]^3$ contains the entire witness support.
For $t\le-7/2$, Cantelli's inequality for the variance-three sum yields
\[
 P(X_1+X_2+X_3\le t)-\Phi(t/\sqrt3)
 \le\frac3{3+t^2}\le\frac{12}{61}<\Lambda\le\Lambda\rho.
\]
For $t\ge3/2$, the discrepancy is at most
\[
 1-\Phi(\sqrt3/2)<\Lambda\le\Lambda\rho.
\]
All these inequalities contradict the choice of the maximizing witness.
The reduction by reflection already handles the other signed discrepancy;
no atom or boundary event is discarded.
\end{proof}

\subsubsection{Exact verification algorithm}\label{ms:finite:n3:sec:check}
Each coefficient is rational. Verification reconstructs the polynomial from its defining features and requires complete threshold coverage. Every unresolved positivity test is rejected.

\paragraph{Piecewise polynomial representation and symmetry}
On each sign orthant, $|y|^3$ is $y^3$ or $-y^3$. Thus $M_c$ is a polynomial
of degree at most five in each variable. Expand this polynomial
from the feature definition, not from a supplied power expansion. Symmetry
permits restriction to $u\le v\le w$. The four initial sign boxes are
\[
 [-7,0]^3,\quad[-7,0]^2\times[0,7],\quad
 [-7,0]\times[0,7]^2,\quad[0,7]^3.
\]
A dyadic box $[L_1,H_1]\times[L_2,H_2]\times[L_3,H_3]$ may be discarded
only if $L_i>H_j$ for some $i<j$. Such a box has no ordered point. All other
accepted inequalities are proved on the whole box, a safe enlargement.

\paragraph{An affine half-space Bernstein test}
For a box, let $B_\alpha$ be the tensor Bernstein basis of degree five in
three variables, so $B_\alpha\ge0$ and $\sum B_\alpha=1$. Write
\[
 M_c=\sum_\alpha b_\alpha B_\alpha,\qquad
 u+v+w-b_j=\sum_\alpha s_\alpha B_\alpha,
\]
where
\[
 s_\alpha=\sum_{i=1}^3\left(L_i+\frac{\alpha_i}{5}(H_i-L_i)\right)-b_j.
\]
\begin{lemma}[Exact affine test]\label{ms:finite:n3:lem:affine}
If a real $\ell\le0$ satisfies $b_\alpha\ge1+\ell s_\alpha$ for all
$\alpha$, then $M_c\ge1$ on the portion $u+v+w\le b_j$.
If $\ell\ge0$ and $b_\alpha\ge\ell s_\alpha$ for all $\alpha$, then
$M_c\ge0$ on the portion $u+v+w\ge b_j$.
\end{lemma}
\begin{proof}
Average the coefficient inequalities with the nonnegative Bernstein weights.
In each case the sign restriction makes $\ell(u+v+w-b_j)$ nonnegative
on the specified half-space.
\end{proof}
Existence of $\ell$ is checked by intersecting rational half-lines. A
coefficient with $s_\alpha=0$ requires $b_\alpha\ge1$ or $b_\alpha\ge0$.
The remaining coefficients give upper or lower rational bounds on $\ell$.
Only signed integer products are needed to compare these bounds.

A box wholly inside the event is checked by $\min b_\alpha\ge1$; a box
strictly outside by $\min b_\alpha\ge0$. A crossing box uses the affine
tests. An unresolved box is bisected along a widest coordinate. Dyadic
subdivision uses the exact degree-five de Casteljau identities
\[
 b_i^{\rm left}=2^{-i}\sum_{j=0}^i\binom ij b_j,
 \qquad
 b_i^{\rm right}=2^{-(5-i)}\sum_{j=i}^5\binom{5-i}{j-i}b_j.
\]
After multiplying the common denominator by $32$, all new numerators are
integers. Ties at the event plane remain on the included side, and that
plane is required to have value at least one.

\paragraph{Moment region and Gaussian enclosure}
The scalar polynomial
\[
 \underline\Phi(a_j/\sqrt3)+\Lambda\rho-q_{c^{(j)}}(\gamma,\rho)
\]
is checked by rational tensor Bernstein coefficients of degree three on
moment rectangles. Here $\underline\Phi$ is a rigorous rational lower bound.
For $s=\gamma^2$ the domain constraint is
\[
 h(s,\rho)=\rho^2(s+4)-(s+2)^2\ge0.
\]
When $\rho\le15/8<2$, this expression decreases in $s\ge0$ and increases
in $\rho\ge0$. If its maximum at the smallest possible $\gamma^2$ and
largest possible $\rho$ on a rectangle is negative, that rectangle is
outside the domain and is discarded. Otherwise Bernstein positivity proves
the scalar inequality or the rectangle is bisected.

For rational $t$ in the required range,
\[
 \Phi(t/\sqrt3)=\frac12+\frac1{\sqrt{6\pi}}
 \sum_{k=0}^{\infty}\frac{(-1)^k t^{2k+1}}{6^k k!(2k+1)}.
\]
Use consecutive alternating partial sums through indices
$34,35$. For $|t|\le4$ the magnitudes decrease, so these enclose the sum.
Rational Machin bounds enclose $\pi$, and integer square roots enclose
$\sqrt{6\pi}$ to rational precision. No floating-point normal cdf enters
acceptance.

\paragraph{The complete finite certificate}
A complete certificate must supply consecutive threshold intervals covering the entire declared range, with 59 rational coefficients for each interval. Acceptance must reconstruct the polynomial expansions and verify the continuous cube and moment inequalities. Undecided cells
are rejected.

The initial expansions are also checked on a complete \(6^3\)
tensor interpolation grid in each sign box. Since the degree is at
most five in each coordinate, agreement on this grid proves the
polynomial identity. This is an independent identity check, not a
second verification of all positivity inequalities.


\subsection{Four summands: a conditional unrestricted bound}\label{ms:finite:n4all:section}

\begin{theorem}[Certificate-conditional four-summand bound]\label{ms:finite:n4all:thm:main}
Assume the three-summand conclusion of Theorem~\ref{ms:finite:n3:thm:main}, as required by the localization below. If the complete bounded-law certificate inequalities specified below
are verified, then for every standardized real law with finite third absolute moment,
\[
 R_4(P)\le\frac{4097}{10000}<\CE.
\]
In particular, $C_4\le4097/10000$. No restriction on the number of
support points is imposed.
\end{theorem}

Put $T=4097/10000$ and $c_0=256/625$. The proof has two steps.
A hypothetical unrestricted maximizer above $T$ lies in a compact
moment class. A finite rational certificate bounds every law in that
class by $T$. The three-point law introduced below is a reference law,
not a reduction of the admissible laws to three points.

\subsubsection{The canonical reference and its excess}
For a standardized $X$, write
\[
 \begin{gathered}
 \vartheta=\E[X|X|],\qquad u=\E X_+=\E X_->0,\qquad
 v_\pm=\frac{1\pm\vartheta}{2},\\
 a=\frac{v_-}{u},\qquad b=\frac{v_+}{u},\qquad
 w=\frac{4u^2}{1-\vartheta^2}.
 \end{gathered}
\]
Both $v_\pm$ are positive. Cauchy's inequality on the two half-lines
gives $0<w\le1$. Define
\begin{equation}\label{ms:finite:n4all:eq:reference}
 Q=\frac{u}{a}\delta_{-a}+(1-w)\delta_0+\frac{u}{b}\delta_b,
 \qquad
 r=\rho(Q)=\frac{1+\vartheta^2}{\sqrt{w(1-\vartheta^2)}}.
\end{equation}
This is a probability law of mean zero and variance one. Its contacts
are
\[
 a=\sqrt{\frac{1-\vartheta}{w(1+\vartheta)}},\qquad
 b=\sqrt{\frac{1+\vartheta}{w(1-\vartheta)}}.
\]

\begin{lemma}\label{ms:finite:n4all:lem:excess}
Set
\[
 W(X)=X_+(X_+-b)^2+X_-(X_--a)^2,\qquad
 \eta=\E W=\rho(P)-r.
\]
Then $\eta\ge0$, and $\eta=0$ if and only if $P=Q$.
Moreover,
\begin{align}
 \E X^3&=b-a+\E[\operatorname{sgn}(X)W(X)],\notag\\
 \E X^4&=a^2-ab+b^2+
   \E[(|X|+2c(X))W(X)],\label{ms:finite:n4all:eq:physical-moments}
\end{align}
whenever the fourth moment is finite, where $c(X)=b$ on $X>0$ and
$c(X)=a$ on $X<0$.
\end{lemma}
\begin{proof}
On a tagged half-line put $y=|X|$ and $c=a$ or $b$. The identities
\[
 y^3=y(y-c)^2+2cy^2-c^2y,
 \qquad
 y^4=(y+2c)y(y-c)^2+3c^2y^2-2c^3y
\]
integrate using $\E[y;\text{tag}]=u$ and
$\E[y^2;\text{tag}]=uc$. They give all three formulas.
If $\eta=0$, the support is contained in $\{-a,0,b\}$; the first
two moments fix its masses as in \eqref{ms:finite:n4all:eq:reference}.
The converse is immediate.
\end{proof}

The excess $\eta$ is not the two-point defect $d$ of Part~I.
Its zero set permits an atom at zero; the reference is two-point
precisely when $w=1$.

\subsubsection{Localization of an unrestricted extremizer}
Write $K_3=59\sqrt3/250$. Theorem~\ref{ms:finite:n3:thm:main} gives
$C_3\le K_3$, and
\[
 c_0^2-K_3^2=\frac{1069}{1562500}>0,\qquad c_0<T.
\]

\begin{lemma}\label{ms:finite:n4all:lem:owner}
If $C_4>c_0$, an unrestricted maximizing pair $(P,H)$ can be chosen
as a positive closed discrepancy. It satisfies
\begin{equation}\label{ms:finite:n4all:eq:owner-class}
 \supp P\subset(-101/25,101/25),\quad
 \rho(P)<46/25,\quad
 \frac{625}{2116}<w\le1,\quad |\vartheta|<\frac{321}{500}.
\end{equation}
\end{lemma}
\begin{proof}
Attainment is Assumption~\ref{ms:finite:inputs}, item~\ref{ms:finite:input:cut49:global:attainment}.
Shevtsova's estimate in Assumption~\ref{ms:finite:inputs}, item~\ref{ms:finite:input:cut49:global:finitecaps}, gives, at this pair,
\[
 \rho<\frac{(3031/10000)(323/500)}{c_0-3031/10000}
       =\frac{979013}{532500}<\frac{46}{25}.
\]
This bound is obtained before any support or fourth-moment assumption.
Put $C=C_4$, $z=H/2$, $F=F_{4,P}(H)$ and $G=F_{3,P}$.
Assumption~\ref{ms:finite:inputs}, item~\ref{ms:finite:input:cut49:global:contact} gives
\begin{align}
 \rho I(y)={}&8\{G(H-y)-F\}+4y\phi(z)+z\phi(z)(y^2-1)\notag\\
 &+\tfrac32C\rho(y^2-1)
   -C\{|y|^3-\rho-3\vartheta y\}.\label{ms:finite:n4all:eq:contact}
\end{align}
Here $I\le0$ everywhere and $I=0$ at every point of the closed
support, including accumulation points. Using the all-law predecessor
bound and $F=\Phi(z)+C\rho/2$, with $s=\sqrt{4/3}$, yields
\begin{align}
 \rho I(y)\le{}&\mathcal H_4(z,y)+4\rho(K_3-C)
                  +4\rho(s-1)K_3\notag\\
 &+C\{-|y|^3+\tfrac32\rho y^2+3\vartheta y-\rho/2\},
 \label{ms:finite:n4all:eq:predecessor}
\end{align}
where
\[
 \mathcal H_4(z,y)=
 8\left[\Phi\!\left(\frac{z-y/2}{\sqrt{3/4}}\right)-\Phi(z)\right]
 +4y\phi(z)+z\phi(z)(y^2-1).
\]

We use the following elementary uniform estimate:
\begin{equation}\label{ms:finite:n4all:eq:gaussian}
 |\mathcal H_4(z,y)|<\frac{\phi(0)}6|y|^3+\frac{|y|}{5}+\frac1{15}.
\end{equation}
Indeed, $\phi(0)<2/5$, $\sup|\phi''|=\phi(0)$ and
$\sup|(x\phi(x))'|=\phi(0)$. Also
\[
 \sup_x|x(3-x^2)\phi(x)|<3/5.
\]
For the last inequality, set $t=x^2$. On $0\le t\le3$,
\[
 \frac94(1+t+t^2/2+t^3/6)-t(3-t)^2
 \ge\frac34\left[7(t-9/14)^2+3/28\right]>0.
\]
Compare with $e^t$ and use $\phi(0)^2<4/25$.
On $t\ge3$, use $|x|^3\phi(x)<1/2$, which follows from $e>8/3$.
For $f(\lambda)=\Phi((z-y/2)/\sqrt{1-\lambda})$,
\[
 f'(0)=\tfrac12(z-y/2)\phi(z-y/2),\qquad
 |f''(\lambda)|<4/15\quad(0\le\lambda\le1/4).
\]
Thus $8\{f(1/4)-f(0)\}=(z-y/2)\phi(z-y/2)+E$, with
$|E|<1/15$. A second-order translation of $\Phi$ gives the cubic
term in \eqref{ms:finite:n4all:eq:gaussian}; the remaining difference
$(z-y/2)\phi(z-y/2)-z\phi(z)$ is bounded by $|y|/5$.

By \eqref{ms:finite:n4all:eq:reference}, $r\le\rho<46/25$ and $w\le1$.
At $|\vartheta|=321/500$,
\[
 \frac{(1+\vartheta^2)^2}{1-\vartheta^2}
   =\frac{124637947681}{36739750000}>(46/25)^2.
\]
Monotonicity in $\vartheta^2$ proves $|\vartheta|<321/500$.
Drop the negative term $4\rho(K_3-C)$ in
\eqref{ms:finite:n4all:eq:predecessor}. Since $4(s-1)<5/8$, its remaining
$C$-bracket is bounded, for $r_y=|y|\ge4$, by
\[
 B(r_y)=-r_y^3+\frac{69}{25}r_y^2
                 +\frac{963}{500}r_y+\frac{23}{100}.
\]
The values of $B,B',B'',B'''$ at 4 are
$-5953/500,-11997/500,-462/25,-6$; hence $B<0$ on that ray.
Only now substitute $C>c_0$ in the negative bracket. Equation
\eqref{ms:finite:n4all:eq:gaussian} gives
\[
 \rho I(y)<P(r_y):=-\frac{643}{1875}r_y^3
 +\frac{17664}{15625}r_y^2+\frac{77257}{78125}r_y
 +\frac{7541}{46875}.
\]
At $101/25$ the values of $P,P',P'',P'''$ are respectively
\[
 -\frac{51557}{9765625},\quad -\frac{520966}{78125},\quad
 -\frac{94558}{15625},\quad -\frac{1286}{625}.
\]
They are negative, so $P<0$ thereafter. Contact excludes every
support point with $|y|\ge101/25$. Finally $r^2\ge1/w$ gives
$w>625/2116$. This proves \eqref{ms:finite:n4all:eq:owner-class}.
\end{proof}

\subsubsection{One defect measure for all frequencies}
For $\eta>0$, let $\nu=W\,dP/\eta$, viewed as a probability on
tagged points $(\sigma,y)$, $\sigma\in\{-1,1\}$, $y>0$.
Write $c_+=b$, $c_-=a$, $p_0=1-w$, and
\[
 k_0(\sigma,y)=\frac1{c_\sigma^2y},\qquad
 k_2(\sigma,y)=\frac1{(y-c_\sigma)^2},\qquad
 k_1(\sigma,y)=\frac1{y-c_\sigma}.
\]
The contacts have zero $\nu$ mass. Singular expressions are integrated
only off the contacts; the following identities prove integrability.

\begin{lemma}\label{ms:finite:n4all:lem:common-measure}
With $g(q)=\varphi_Q(q)$ and
$r_q(\sigma,y)=\exp(\mathrm i\sigma q\,\cdot)[0,c_\sigma,c_\sigma,y]$,
\begin{align}
 \varphi_P(q)-g(q)&=\eta\int r_q\,d\nu,\notag\\
 \eta\int k_0\,d\nu&=p_0-P\{0\}\le p_0,\notag\\
 \eta\int_{\sigma}k_2\,d\nu&=u-c_\sigma P\{\sigma c_\sigma\}\le u,
 &\int_{\sigma}k_1\,d\nu&=0.\label{ms:finite:n4all:eq:budgets}
\end{align}
The same $\nu$ occurs at every frequency, and $|r_q|\le |q|^3/6$.
For $\eta=0$, $P=Q$ and every characteristic-function difference is
zero, without defining $\nu$.
\end{lemma}
\begin{proof}
Hermite interpolation at $0,c,c$ has remainder
$y(y-c)^2\exp(\mathrm i\sigma q\,\cdot)[0,c,c,y]$.
On each tag, the linear and quadratic terms have the same integrals
against $P$ and $Q$, by their matching first two tagged moments.
Both interpolating polynomials have constant term 1. Their constant
contributions agree only after summing the tags and adding the atom
at zero, since both total masses are 1.
Hermite--Genocchi gives the response bound. Multiplication of the
three costs by $W=y(y-c)^2$ gives
$(y-c)^2/c^2$, $y$, and $y(y-c)$, respectively.
Their integrals give \eqref{ms:finite:n4all:eq:budgets}, including the omitted
contact masses. Absolute integrability of $k_1$ follows from Cauchy
and the $k_2$ bound. Lemma~\ref{ms:finite:n4all:lem:excess} handles $\eta=0$.
\end{proof}

Fix a threshold $H$ and finite coefficients $\kappa_j(H)$, defined
below. For the bounded class $|X|\le R_*=17/4$, all suprema below
are over $\sigma\in\{-1,1\}$ and $0<y\le R_*$, off the contact.
The linear response is
\[
 A_H(\sigma,y)=\Re\sum_{j=0}^{N-1}
        4\kappa_j(H)g(q_j)^3r_{q_j}(\sigma,y).
\]
For any fixed $\lambda_0,\lambda_+,\lambda_-\ge0$ and
$\theta_+,\theta_-\in\R$, a bound
\begin{equation}\label{ms:finite:n4all:eq:priced-supremum}
 S_H\ge\sup_{\sigma,y}
 \{A_H-\lambda_0k_0-\lambda_\sigma k_2-\theta_\sigma k_1\}
\end{equation}
implies
\begin{equation}\label{ms:finite:n4all:eq:gamma}
 \eta\int A_H\,d\nu\le
 \Gamma_H(\eta):=p_0\lambda_0+u(\lambda_++\lambda_-)+S_H\eta.
\end{equation}
Upper bounds for $p_0,u$ may replace them. The two equality prices
retain their signs. In particular, $S_H$ and $\Gamma_H$ need not be
nonnegative. The prices are fixed before taking the supremum over
the whole tagged domain. This preserves the cancellation between
frequencies; separate absolute estimates for each $r_q$ do not.

\subsubsection{The fourth-power remainder}
For complex $\psi=g+\delta$,
\begin{equation}\label{ms:finite:n4all:eq:quartic}
 \psi^4-g^4=4g^3\delta+
 \delta^2(\psi^2+2\psi g+3g^2).
\end{equation}
Suppose $|g|\le A\le1$, $|\psi|\le1$ and $|\delta|\le E$.
Two complete bounds for the remainder are
\begin{equation}\label{ms:finite:n4all:eq:quartic-basic}
 E^2(p^2+2pA+3A^2),\quad p=\min(1,A+E),\qquad
 v(4A^2+2-v),\quad v=\min(E^2,2A^2+1).
\end{equation}
The first follows directly from \eqref{ms:finite:n4all:eq:quartic}. For the
second, use
\[
 \psi^4-g^4-4g^3\delta
   =12\delta^2\int_0^1(1-t)(g+t\delta)^2\,dt
\]
and
$|g+t\delta|^2=(1-t)|g|^2+t|\psi|^2-t(1-t)|\delta|^2$.

A sharper two-disk bound is also finite algebra. Rotate $g=a\ge0$
and put $\delta=x+\mathrm i y$, $s=|\delta|^2$. Then
\[
 |\psi^4-g^4-4g^3\delta|^2
 =s^2\{36a^4+4a^2s+s^2+8ax(6a^2+s)+24a^2x^2\}.
\]
For fixed $a,s$, the expression in braces is convex in $x$.
Its maximum on the two disks occurs at a collinear endpoint or on
$|\psi|=1$. Put $r_+=\min(E,1)$, $m=\min(A,1-r_+)$,
$r_-=\min(E,1+A)$. The collinear bounds are
\[
 r_+^2(6m^2+4mr_++r_+^2),\qquad
 r_-^2(6A^2-4Ar_-+r_-^2),\qquad
 \begin{cases}r_-^4,&r_-\le1,\\
 r_-^2(3r_-^2-8r_-+6),&r_-\ge1.
 \end{cases}
\]
On the unit circle it remains to maximize
\begin{equation}\label{ms:finite:n4all:eq:lens-polynomial}
 Q_A(s)=s^2\{18A^4+12A^2+6-(12A^2+8)s+3s^2\}
\end{equation}
on $[(1-A)^2,\min(E^2,(1+A)^2)]$. Its endpoints and the sole
interior local maximum
\[
 s_-=(3A^2+2-A\sqrt{4-3A^2})/2
\]
suffice. For $A=0$ the interval is degenerate or empty.
Taking the maximum of the collinear values and the square root of
this last maximum gives the lens bound. To see that no $a<A$ is
lost, note that the unit-circle expression is convex in $a^2$;
its lower feasible endpoint is already collinear. The collinear
expressions increase with their admissible radius after maximization
over $a$. Thus these are whole-disk bounds, not sampled values.
Their minimum with \eqref{ms:finite:n4all:eq:quartic-basic} is valid.

On a closed excess interval $[L,U]$, an affine radius
$e_j(\eta)=A_j+B_j\eta$ is used only when it is nonnegative at both
ends and bounds $|\delta_j|$ throughout. Freeze
$p_j=\min(1,a_j+\max\{e_j(L),e_j(U)\})$ in the first bound of
\eqref{ms:finite:n4all:eq:quartic-basic}. This gives a quadratic in $\eta$
with nonnegative leading coefficient. Alternatively, freezing $E$
on the whole cell gives a constant disk or lens bound. A choice
between complete bounds is allowed; a coefficientwise minimum is not.

\subsubsection{A physical Prawitz comparison}
We use the closed-CDF Prawitz majorant~\cite{Prawitz1972}, in the
normalization of Assumption~\ref{ms:finite:inputs}, item~\ref{ms:finite:input:orig:lem:BPcardinal}. To distinguish the
factor from the kernel $K$ used there, put
\[
 \widehat K(u)=2K(u)=(1-u)+\mathrm i R(u),\qquad
 R(u)=(1-u)\cot(\pi u)+1/\pi.
\]
For $Z=(X_1+\cdots+X_4)/2$,
\begin{equation}\label{ms:finite:n4all:eq:prawitz}
 F_Z(x)\le J_P(x):=\frac12+
 \int_0^1\Re\{\widehat K(u)e^{-6\mathrm i ux}\varphi_P(3u)^4\}\,du.
\end{equation}
The real integrand is regular at zero; no convergence of the complex
integral by itself is asserted. With $D=x-Z$ and $S(u)=uR(u)$ it is
\[
 F_x(u)=(1-u)\E\cos(6Du)+S(u)\E[6D\operatorname{sinc}(6Du)].
\]
The cotangent partial fractions, extended to the endpoints, give
\[
 0\le S\le1/3,\qquad |S'|<2,\qquad |S''|<5.
\]
Here is a direct derivative check. Separate the first pole:
\[
 S=\frac1\pi-\frac2\pi(A+MG),\quad
 A=\frac{u^2}{1+u},\quad M=u^2(1-u),\quad
 G=\sum_{j\ge2}\frac1{j^2-u^2}.
\]
On $[0,1]$, $|A'|\le3/4$, $|A''|\le2$, $|M|\le1/6$,
$|M'|\le1$ and $|M''|\le4$. The derivative series converge
uniformly, and
\[
 G\le3/4,\qquad |G'|\le10/27<1/2,\qquad
 |G''|\le106/135<7/6.
\]
The first sum telescopes. For the other two use
$j^2-u^2\ge3j^2/4$ and
$\sum_{j\ge2}j^{-k}\le2^{-k}+\int_2^\infty t^{-k}\,dt$,
with $k=4,6$. The product rule gives
$|S'|<19/18<2$ and $|S''|<223/54<5$.
For $u>1/2$, put $v=1-u$ and use
$v\cot(\pi v)\le1/\pi$; for $u\le1/2$ the cotangent is
nonnegative. This proves $R\ge0$. The positive partial-fraction
terms then give $0\le S\le1/\pi<1/3$.
Since $\operatorname{sinc}(v)=\int_0^1\cos(tv)\,dt$,
\begin{equation}\label{ms:finite:n4all:eq:integrand-curvature}
 |F_x''(u)|\le42\E|D|+108\E D^2+24\E|D|^3.
\end{equation}
These differentiations are justified by the displayed moments.

Now assume $|X|\le R_*=17/4$ and $\rho<46/25$.
The fourth moment follows from these hypotheses; it is not an
additional input. Write $\kappa=\E X^3$, $M_4=\E X^4$.
Lemma~\ref{ms:finite:n4all:lem:excess} gives the useful cell bounds
\begin{align}
 |\kappa|&\le\min\{\rho,|b-a|+\eta\},\notag\\
 M_4&\le\min\{R_*\rho, a^2-ab+b^2+
                         (R_*+2\max(a,b))\eta\}.
 \label{ms:finite:n4all:eq:local-moments}
\end{align}
For $x=H/2$, choose whole-cell bounds $\kappa_*$ and $M_{4,*}$
from these formulas and put
\[
 D_2=x^2+1,\qquad
 D_4=x^4+6x^2+2|x|\kappa_*+(M_{4,*}+9)/4.
\]
Then $\E|D|\le\sqrt{D_2}$ and
$\E|D|^3\le\sqrt{D_2D_4}$. Uniformly on $|x|\le2$, the simpler
bounds $\E D^2\le5$, $\E D^4<54$, $\E|D|<9/4$,
$\E|D|^3<33/2$ give $|F_x''|<2061/2$.

Take $N=256$, $u_j=(j+1/2)/N$, $q_j=3u_j$ and
\[
 \kappa_j(H)=\widehat K(u_j)e^{-\mathrm i q_jH}/N,\qquad
 C_H(Q)=1+2\Re\sum_j\kappa_j(H)g(q_j)^4-2\Phi(H/2).
\]
Composite midpoint quadrature gives one raw numerator reserve
\begin{equation}\label{ms:finite:n4all:eq:midpoint}
 Q_H=\min\left\{\frac{687}{8N^2},
 \frac{42\sqrt{D_2}+108D_2+24\sqrt{D_2D_4}}{12N^2}\right\}.
\end{equation}
It controls the integral of the original law $P$ once. No separate
reference, linear-transfer or remainder quadrature error is added.

The raw thresholds are $H_\ell=-4+\ell/8$, $0\le\ell\le57$.
To fill the intervals between them, bound the characteristic function
on entire frequency cells, not just at the midpoints. If
$a_j\ge|g(q_j)|$ and $e_j$ bounds $|\varphi_P(q_j)-g(q_j)|$ on the
whole excess cell, then
\[
 P_j=\min\{1,a_j+e_j+3/(2N)\}
\]
bounds $|\varphi_P(3u)|$ throughout $j/N\le u\le(j+1)/N$,
since $\E|X|\le1$. Two complete kernel bounds are
\[
 u^2|\widehat K(u)|\le u^2(1-u)+u/3,\qquad
 u^2|\widehat K(u)|\le u^2(1-u)+u(1-u)/2.
\]
The second follows by retaining the first term of the cotangent
partial fractions:
$S(u)\le(1-u)(1+2u)/(\pi(1+u))<(1-u)/2$.
Let $d_j$ be the smaller of the two exact integrals over the $j$th
frequency cell. Their separate totals are $1/4$ and $1/6$.
Thus, for the same whole geometry and excess cell,
\[
 |J_P''(x)|\le B=36\sum_jd_jP_j^4.
\]
Since $\sup|\Phi''|<1/4$, the smooth physical function
$H\mapsto2\{J_P(H/2)-\Phi(H/2)\}$ needs the single interpolation
reserve
\begin{equation}\label{ms:finite:n4all:eq:H-reserve}
 E_H=\frac{h^2}{16}(B+1/4),\qquad h=1/8.
\end{equation}
The elementary cube, mass and unit radii
\[
 \eta q_j^3/6,\qquad \max(a,b)^2p_0q_j^2,\qquad 1+a_j
\]
suffice for this curvature calculation. The mass bound follows from
$|r_q|\le q^2/y$ and the first budget in
\eqref{ms:finite:n4all:eq:budgets}.

For $H\le-4$, Cantelli gives $F_Z(H/2)\le1/5$.
For $H\ge25/8$, $\Phi(H/2)>4/5$. The latter inequality follows,
for example, by integrating the alternating Gaussian Taylor
polynomial through degree six on $[0,25/16]$ and using
$\phi(0)>19947/50000$. Consequently both threshold tails have
positive normalized score at most $2/5$.

\subsubsection{The low-excess certificate}
On a geometry rectangle and a closed excess cell $[L,U]$, let
$M_H$ bound $C_H(Q)-Tr$ for the entire rectangle. Let
$\Gamma_H$ be obtained from \eqref{ms:finite:n4all:eq:gamma}, and let
\[
 R_2(\eta)=r_0+r_1\eta+r_2\eta^2,\qquad r_2\ge0,
\]
bound the sum of the fourth-power remainders weighted by
$|\kappa_j(H)|$. Combining \eqref{ms:finite:n4all:eq:quartic} at the same
$H$ gives
\begin{equation}\label{ms:finite:n4all:eq:low-polynomial}
 2\{F_{4,P}(H)-\Phi(H/2)\}-T\rho
 \le P_H(\eta):=M_H+Q_H+E_H+2\Gamma_H(\eta)+2R_2(\eta)-T\eta.
\end{equation}
When used at grid points, $E_H$ reserves interpolation to adjacent
thresholds. The leading coefficient is nonnegative. Endpoint
nonpositivity therefore proves the bound on the entire closed excess
cell. Equivalently its degree-two Bernstein coefficients are
\[
 P_H(L),\quad
 \frac{P_H(L)+P_H(U)}2-\frac{p_2(U-L)^2}{2},\quad P_H(U),
\]
where $p_2$ is the quadratic coefficient. For a quantitative signed
score, with $r_L\le r$ and $\rho<46/25$, use
\[
 T+\max_{\eta=L,U}
 \left\{\frac{P_H(\eta)}{46/25},
        \frac{P_H(\eta)}{r_L+\eta}\right\}.
\]
Both denominator faces are needed when the numerator is negative.
The cap branch is convex; for the other branch the derivative
numerator has derivative $P_H''(\eta)(r_L+\eta)\ge0$, so an
interior critical point can only be a minimum.

Here are the finite enclosures that produce these polynomials.
The center is enclosed by signed Taylor expansion in $(w,\vartheta)$
or in the contacts $(a,b)$ along probability-preserving paths.
For either pair of coordinates $i,j$,
\[
 \partial_i C_H=8\Re\sum_j\kappa_jg^3\partial_i g,\qquad
 |\partial_{ik}C_H|\le8\sum_j|\kappa_j|
       (3D_iD_k+D_{ik}),
\]
where $D_i,D_{ik}$ bound the derivatives of the explicit
three-atom characteristic function. The first-order terms retain
their signs; the second-order error bounds the whole rectangle.
The function $r(w,\vartheta)$ is jointly convex. Indeed, writing
$r=w^{-1/2}f(\vartheta)$, $f=(1+\vartheta^2)/\sqrt{1-\vartheta^2}$,
the Hessian determinant has the positive factor
\[
 3ff''-(f')^2=
 \frac{9+9\vartheta^2+15\vartheta^4-\vartheta^6}{(1-\vartheta^2)^3}.
\]
Subtracting its supporting plane therefore encloses $C_H-Tr$
without an absolute Hessian charge for $r$.

For the signed response, put $R_*=17/4$, $y_k=k/64$,
$0\le k\le272$, and use their nearest-node closed cells,
clipped to $[0,R_*]$. Hermite--Genocchi gives
\begin{align}
 r_q(\sigma,y)&=\sum_{m=0}^{60}(\mathrm i\sigma)^{m+3}b_m h_m+E_q,
 &|E_q|&\le\frac{q^3}{6}\frac{(qR_*)^{61}}{61!},\notag\\
 b_0&=q^3/6,&b_m&=b_{m-1}qR_* /(m+3),\notag\\
 h_0&=1,&h_m&=(y/R_*)h_{m-1}+(m+1)(c/R_*)^m.
 \label{ms:finite:n4all:eq:response-series}
\end{align}
The bound $|\partial_y r_q|\le q^4/24$ adds the half-cell reserve
$q^4/(48\cdot64)$. The signed frequency coefficients are summed
before bounding their polynomial on a $y$ cell. If $g_0,c_0'$ are
nominal values, $e_q\ge|g-g_0|$ is the geometry error,
$p_q\ge|g|$ and $a_{0,q}\ge|g_0|$, then
\[
 |g^3-g_0^3|\le e_q(p_q^2+p_qa_{0,q}+a_{0,q}^2),
 \qquad |\partial_c r_q|\le q^4/12.
\]
These give a whole-rectangle correction for
$g^3r_q-g_0^3r_{q,0}$, with no second copy of the series or mesh
reserve. Thus the priced supremum bounds actual references throughout
the rectangle, not only its nominal center.

The prices used in the finite certificate are chosen from
\begin{align*}
 \lambda_0&\in\{0,1/1000,1/400,1/200,1/100,1/50,1/25,
                         2/25,4/25,8/25,16/25\},\\
 \lambda_\pm&\in\{0,1/2000,1/1000,1/400,1/200,1/100,1/50\},\\
 \theta_\pm&\in\{0,\pm1/200,\pm1/100,\pm1/50\}.
\end{align*}
For each fixed choice, all 273 cells on each of the two tags enter
the supremum.
When $\lambda_\sigma=0$, only $\theta_\sigma=0$ is used;
otherwise the reciprocal contact rays would give an infinite supremum.
On a cell $[y_L,y_U]$, with $c\in[c_L,c_U]$,
$k_0\ge1/(c_U^2y_U)$. The remaining cost is the concave quadratic
$-\lambda_\sigma z^2-\theta_\sigma z$ in $z=1/(y-c)$.
Its endpoints and admissible vertex suffice; when the cell crosses
the contact both reciprocal rays are included. No division at the
contact or at zero is used for the response series.

The same budgets yield whole-cell radius models by pricing
$|r_q|$ or its signed coordinate projections. They are combined
with the three elementary radii above and the complete two-disk
remainder. An independent absolute branch uses
\[
 |\psi^4-g^4|\le E(p^3+p^2A+pA^2+A^3),\qquad p=\min(1,A+E),
\]
or the complete linear-plus-remainder bound. It is evaluated as a
whole alternative. In every branch, the actual reference modulus,
the actual normalization $\rho=r+\eta$, one midpoint reserve and
one physical threshold reserve are retained. The final maximum over
$H$ is taken only after assembling each complete row.

\subsubsection{Closing the large excess}
Let $p=P(X\ge0)$ and $q=1-p$; zero belongs to the nonnegative class.
Condition on exactly $j$ nonnegative observations among four.
For $Z=(X_1+\cdots+X_4)/2$, the conditional moments are
\[
 \mu_j=\frac{u(j-4p)}{2pq},\qquad
 V_j=\frac14\left[j\left(\frac{v_+}{p}-\frac{u^2}{p^2}\right)
 +(4-j)\left(\frac{v_-}{q}-\frac{u^2}{q^2}\right)\right].
\]
Closed-event Cantelli gives, for every $x$,
\begin{equation}\label{ms:finite:n4all:eq:count-cantelli}
 F_Z(x)\le\sum_{j=0}^4\binom4j p^jq^{4-j}A_j(x),\qquad
 A_j(x)=\begin{cases}1,&x\ge\mu_j,\\
 V_j/(V_j+(\mu_j-x)^2),&x<\mu_j.
 \end{cases}
\end{equation}
This includes variance zero and atoms at the threshold. The necessary
sign-mass interval is
\[
 w(1-\vartheta)/2\le p\le1-w(1+\vartheta)/2.
\]
It is covered by closed cells. On a threshold cell $[x_L,x_U]$,
bound $A_j(x_U)$ and subtract a lower bound for $\Phi(x_L)$.
The five degree-four Bernstein coefficients bound the binomial
mixture jointly; its masses are not maximized independently.
Each retained sign-mass cell supplies a nonnegative raw numerator
bound $d_i$ and a correlated lower bound $r_i\le r$, so its score
is at most $d_i/(r_i+\eta)$.

If $r_B$ is the whole-box lower bound for $r$, put
\[
 c_C^0=\max\left\{0,\max_i(d_i/T-r_i),\frac{2}{5T}-r_B\right\},
 \qquad
 c_C=10^{-4}\lceil10^4c_C^0\rceil+10^{-4}.
\]
This closes every $\eta\ge c_C$ by a strict bound below $T$,
including the threshold tails. Independently, Shevtsova closes
every $\rho\ge979013/533000$:
\[
 \frac{3031}{10000}\left(1+\frac{323}{500\rho}\right)\le T.
\]
Hence $c_S=\max\{0,979013/533000-r_B\}$ is another closed sufficient
cutoff. Take $c_B=\min(c_C,c_S)$, selecting one complete route.
The low-excess interval $[0,c_B]$ and the large-excess range
$[c_B,\infty)$ include their common endpoint.

\subsubsection{The finite cover and conclusion}
\begin{proposition}[Bounded-law certificate implication]
\label{ms:finite:n4all:prop:certificate}
Assume all the terminal inequalities and domain covers in the following
finite scheme have been verified. Then, for every standardized law $P$ with
\[
 |X|\le17/4,\quad \rho(P)<46/25,\quad
 625/2116\le w\le1,\quad |\vartheta|\le13/20,
\]
and every real $H$,
\[
 \frac{2\{F_{4,P}(H)-\Phi(H/2)\}}{\rho(P)}\le T.
\]
The same assertion holds for the genuine reflected law $-P$.
\end{proposition}
\begin{proof}
Divide $[625/2116,1]$ into 20 equal closed bands and $[0,13/20]$
into 13 closed bands of width $1/20$. There are 260 rectangles.
The minimum of
$r^2=(1+\vartheta^2)^2/[w(1-\vartheta^2)]$ is attained at the
largest $w$ and smallest $|\vartheta|$. A rectangle with
$\min r^2\ge(46/25)^2$ is empty for the strict physical cap.
For every other rectangle construct its own large-excess cutoff.
If $c_B=0$, the large-excess bound covers all $\eta\ge0$ directly.
If $c_B>0$, divide $[0,c_B]$ into 16 equal closed intervals. At each interval
check all 58 thresholds, their 57 joins, the two tails, and both
actual reflected orientations, using the preceding inequalities.
A parent is accepted by its complete bound or by complete bounds
for all four children obtained by bisecting both coordinates.
Each child has its own cutoff; no part of another child's domain
or numerical conclusion is used in its place.

The certificate condition requires a disposition for every one of the
260 closed parent rectangles: a proved necessary-condition contradiction,
a complete direct comparison, or complete comparisons on all four children.
Acceptance requires the complete rational inequalities and their
reconstruction, with no unresolved rectangle.
All elementary operations
are outward rational enclosures; square roots are checked by integer
squaring. Trigonometric, Gaussian and cotangent values are enclosed
by finite series with analytic remainders, using the arithmetic
language of Assumption~\ref{ms:finite:inputs}, item~\ref{ms:finite:input:cut49:short:arithmetic}. The response uses
\eqref{ms:finite:n4all:eq:response-series}, not divided quotients at contacts.
The assumed complete arithmetic reconstruction must reproduce the whole cover
and its upper bounds; no sampled value or optimizer's status supplies
an inequality.

For completeness, acceptance has a direct mathematical interpretation.
Every actual canonical point belongs to a closed parent.
In a nonempty parent or child, the large-excess inequalities cover
$\eta\ge c_B$, while the low-excess inequalities cover all of
$[0,c_B]$, including $\eta=0$. The same unknown law supplies the
single $\nu$ in every priced frequency response. Quadrature and
physical interpolation extend the finite rows to every central
threshold; the tails cover the rest. Thus the accepted rational
inequalities prove the proposition for all actual laws, including
bounded infinite supports and zero atoms.
\end{proof}

\begin{proof}[Proof of Theorem~\ref{ms:finite:n4all:thm:main}]
Suppose $C_4>T$. Since $T>c_0$, Lemma~\ref{ms:finite:n4all:lem:owner} gives
an unrestricted positive closed maximizing pair in the class of
Proposition~\ref{ms:finite:n4all:prop:certificate}. The proposition bounds its
score by $T$, a contradiction. Both discrepancy signs are covered:
\[
 \Phi(x)-F_Z(x)\le\Phi(x)-F_Z(x-)
                     =F_{-Z}(-x)-\Phi(-x).
\]
This is the closed-CDF reflection identity, with the atom retained.
Finally, elementary Machin bounds give $\pi<3927/1250$ and hence
$\phi(0)>19947/50000$; also $\sqrt{10}>15811/5000$. Therefore
\[
 \CE>\frac{204862339}{500000000}
       =\frac{4097}{10000}+\frac{12339}{500000000}>T.
\]
No equality case $C_4=T$ is excluded.
\end{proof}


\subsection{Five summands: the conditional composition}
\label{ms:finite:five:section}

This section specializes the common contact, Fourier, resource and count
arguments to five summands and proves a separate support-free comparison
for the remaining moment corner. Its conclusion is conditional on the finite partitions and numerical
enclosures specified below.

\subsubsection{The contrary witness and its predecessor allowance}

\begin{lemma}[Zero allowance at the fifth index]\label{ms:finite:five:lem:qzero}
Assume $C_4\le\CE$ and $C_5>\CE$. There is an attaining standardized
marginal $P$ and a finite raw threshold $H$ with positive closed-CDF
discrepancy. For that same marginal,
\begin{equation}\label{ms:finite:five:eq:qzero}
 q_5=5\bigl(R_4(P)-C_5\bigr)_+=0.
\end{equation}
Its support is compact before any higher moments are invoked.
\end{lemma}
\begin{proof}
Assumption~\ref{ms:finite:inputs}, item~\ref{ms:finite:input:cut49:global:attainment} supplies the attained positive closed
edge, with reflection of the whole marginal when required. Since
$R_4(P)\le C_4\le\CE<C_5$, \eqref{ms:finite:five:eq:qzero} follows.
Assumption~\ref{ms:finite:inputs}, item~\ref{ms:finite:input:cut49:global:contact} then gives compact support from contact
with a positive cubic leading coefficient. No fourth or fifth moment is
used in either step. In particular, the zero allowance is not an
assertion that $C_5\le C_4$.
\end{proof}

The quantitative support argument of
Assumption~\ref{ms:finite:inputs}, item~\ref{ms:finite:input:cut49:support:theorem} is applied at the actual index $n=5$
with \eqref{ms:finite:five:eq:qzero}. A finite support certificate must prove its
pointwise contact exclusion on both exterior support rays and the entire
threshold line, including both infinite threshold tails. A finite mesh
without its whole-cell remainder does not establish this implication.
Only after this acquisition may the supported moment and Fourier
inequalities in Assumption~\ref{ms:finite:inputs}, items
\ref{ms:finite:input:afcore:sec:fourier},
\ref{ms:finite:input:afcore:sec:conserved} and
\ref{ms:finite:input:rescore:sec:setup}, be used.

\subsubsection{The complete moment-domain cover}\label{ms:finite:five:sec:cover}

Write $\beta=\E[X|X|]$ and $m=\E X_+=\E X_->0$.
Lemma~\ref{ms:sign-moment} gives $|\beta|<1$ and the necessary restriction
\begin{equation}\label{ms:finite:five:eq:beta}
 \rho\ge\frac{1+\beta^2}{\sqrt{1-\beta^2}}.
\end{equation}
Define
\[
 b(r)=\left(\frac{r\sqrt{r^2+8}-r^2-2}{2}\right)^{1/2}.
\]
Solving \eqref{ms:finite:five:eq:beta} shows that $\rho\le r$ implies
$|\beta|\le b(r)$. The finite specification uses
\begin{equation}\label{ms:finite:five:eq:benclosures}
 b_2=\frac{84729828287999}{250000000000000},\qquad
 b_3=\frac{35139006743001}{10^{14}}.
\end{equation}
Before using these values the scalar certificate must verify
$b_2\ge b(237/200)$ and $b_3\ge b(6/5)$, with positive radicals
and outward rounding. These are finite exact inequalities, not measured
approximations to a moment domain.

The proposed low-moment domain decomposition is the following. In the
table, ``full'' means the whole necessary signed beta range from
\eqref{ms:finite:five:eq:beta}. Every interval is closed. The gapped method means
resource acquisition followed, if needed, by the count comparison; it
does not assume a support gap at the beginning of the row.
\begin{center}\small
\begin{tabular}{@{}llll@{}}\toprule
Region & $\rho$ interval & $\beta$ interval & Method\\\midrule
B0 & $[1,11/10]$ & full & resources/counts\\
BN & $[11/10,23/20]$ & full & resources/counts\\
B2 S0 & $[23/20,467/400]$ & $[-b_2,1/5]$ & resources/counts\\
B2 S1 & $[23/20,467/400]$ & $[1/5,b_2]$ & resources/counts\\
B2 S2 & $[467/400,237/200]$ & $[-b_2,1/5]$ & resources/counts\\
B2 S3 & $[467/400,237/200]$ & $[1/5,b_2]$ & resources/counts\\
B3 S0 & $[237/200,477/400]$ & $[-b_3,1/5]$ & resources/counts\\
B3 T0 & $[237/200,477/400]$ & $[1/5,3/10]$ & resources/counts\\
B3 T1 & $[237/200,477/400]$ & $[3/10,b_3]$ & resources/counts\\
B3 S2 & $[477/400,6/5]$ & $[-b_3,1/5]$ & resources/counts\\
B3 T2 & $[477/400,6/5]$ & $[1/5,3/10]$ & resources/counts\\
B3 lower & $[477/400,6/5]$ & $[3/10,17/50]$ & resources/counts\\
B3 corner & $[477/400,6/5]$ & $[17/50,b_3]$ & gap-free\\\bottomrule
\end{tabular}
\end{center}
This is a domain cover once the scalar enclosures just stated are proved.
Each gapped row requires a complete closed cover of $p\in[0,1]$;
the impossible faces $p=0,1$ are excluded by centering and unit variance.
On every other panel either a proved resource contradiction excludes the
law or a simultaneously acquired contract passes to all six original
nearest-count roots. Signed beta scopes are not identified by reflection.
In particular, the boundary $\beta=17/50$ belongs to both adjacent pieces.

The proposed supported Fourier domains above this cover are
\begin{center}
\begin{tabular}{@{}lc@{}}\toprule
$\rho$ interval & Required acquired support radius $R$\\\midrule
$[6/5,5/4]$ & $13/5$\\
$[5/4,7/5]$ & $117/40$\\
$[7/5,3/2]$ & $25/8$\\
$[3/2,8/5]$ & $133/40$\\
$[8/5,17/10]$ & $7/2$\\
$[17/10,7361/4000]$ & $15/4$\\\bottomrule
\end{tabular}
\end{center}
For each row, the support certificate must first acquire $|X|\le R$.
A valid actual-index route is Lemma~\ref{ms:legacy-support:direct}, with
$n=5$, zero predecessor allowance, and
\[
 S=1119/1000,\qquad E=91/1000,\qquad c=819/2000.
\]
For the six rows in their displayed order, take $r$ to be the upper
$\rho$ endpoint and take the respective upper bounds on $|\beta|$ to be
\[
 0.389174,\quad 0.478710,\quad 0.525521,\quad 0.565414,\quad
 0.600001,\quad 0.641563.
\]
They follow from the signed-second-moment inequality, since in each row
$(1+b^2)^2>r^2(1-b^2)$ and $0<b<1$. Exact substitution gives
$Q(R),Q'(R),Q''(R),L(R),L'(R)<0$ for that lemma's polynomials;
$S^2>5/4$ and $((E+11/2)/5)^2>5/4$ give the scale bounds.
The remaining support conditions are precisely the six pairs of
whole-threshold endpoint inequalities
\[
 \sup_{H\in\R}\mathcal H_5(H,\eta R)<-cQ(R),\qquad \eta\in\{-1,1\}.
\]
Those inequalities are not proved by these rational ray gates. The
complete-threshold method in Appendix~\ref{ms:legacy-support:section}
specifies how a finite certificate could establish them, including both
infinite tails. The newer support theorem for $n\ge42$ is not applicable
at this step.

The Fourier certificate must then majorize all four terms of
the Prawitz bound in Assumption~\ref{ms:finite:inputs}, item~\ref{ms:finite:input:afcore:eq:prawitz}, including endpoint pieces and the infinite
normal tail, on a complete closed strip cover. Its required conclusion
is an upper score below $c_0:=4097/10000$.

For the remaining moment tail the published iid structural estimate
\cite{Shevtsova2013} gives
\[
 R_5(P)\le\frac{3031}{10000}
             \left(1+\frac{323}{500\rho}\right).
\]
The scalar certificate must check its value at $\rho=7361/4000$ is
at most $819/2000$ and that $819/2000<c_0<\CE$.
The displayed majorant decreases in $\rho$, so this single endpoint
comparison covers the entire infinite tail. The common smoothing and
constant-enclosure arguments are required by Assumption~\ref{ms:finite:inputs},
items~\ref{ms:finite:input:elementary:section} and~\ref{ms:finite:input:afcore:sec:fourier}.

\subsection{A direct sign-count inequality without a support gap}
\label{ms:finite:five:sec:gapfree}

The inequalities in this section apply to every standardized common iid
marginal in their stated domain. They require neither extremality nor
contact, bounded support, a predecessor estimate, or a physical support gap.
Put $p=\Pp(X>0)$, $q=1-p$, $v=pq$ and $h=v/m$.
Both $p,q$ are positive. Assign zero atoms to the nonpositive class and set
\[
 B=\mathbf 1_{\{X>0\}},\qquad Y=p+hX,\qquad W=Y-B.
\]
Then $\E(W\mid B)=0$. Direct conditional second-moment calculation gives
\begin{align}
 v_0&=p^2\left(\frac{q(1-\beta)}{2m^2}-1\right),
       \label{ms:finite:five:eq:v0}\\
 v_1&=q^2\left(\frac{p(1+\beta)}{2m^2}-1\right).
       \label{ms:finite:five:eq:v1}
\end{align}
For $J=B_1+\cdots+B_5$, write
\[
 b_j(p)=\binom5j p^jq^{5-j},\qquad
 V_j=(5-j)v_0+jv_1.
\]
Conditioning first on a particular sign vector gives independent
conditional variables. Every sign vector with count $j$ gives the same
sum law, of variance $V_j$, which justifies conditioning only on $J=j$.
This is the specialization of the common representation in
\eqref{ms:finite:countcore:representation}.

\begin{lemma}[Gap-free count bound]\label{ms:finite:five:lem:cantelli}
Define
\[
 T(V,d)=\begin{cases}1,&d\le0,\\ V/(V+d^2),&d>0.\end{cases}
\]
For every real $z$,
\begin{equation}\label{ms:finite:five:eq:count}
 \Pp\Bigl(\sum_{i=1}^5Y_i\le z\Bigr)
 \le Q(z):=\min\left\{1,\sum_{j=0}^5b_j(p)T(V_j,j-z)\right\}.
\end{equation}
At $H=(z-5p)/h$, the normalized positive closed-CDF error is at most
\begin{equation}\label{ms:finite:five:eq:score}
 \frac{\sqrt5}{\rho}
 \left[Q(z)-\Phi\left(\frac{(z-5p)m}{\sqrt5\,pq}\right)\right].
\end{equation}
\end{lemma}
\begin{proof}
For a centered variable $U$ of variance $V$ and $d>0$, Markov applied to
$(U-V/d)^2$ on $\{U\le-d\}$ gives
$\Pp(U\le-d)\le V/(V+d^2)$. If $V=0$, that event is empty.
At $d=0$ the safe bound is one even when $V=0$. Apply this inequality
separately to each conditional sum and use its actual binomial mass.
Finally $\sum Y_i=5p+hS_5$ gives the Gaussian coordinate in
\eqref{ms:finite:five:eq:score}.
\end{proof}
In particular, $T(0,0)=1$ retains a possible atom on an integer count
boundary. It cannot be replaced by zero.

\subsubsection{Whole-box bounds in the remaining corner}

The additional necessary moment restrictions are
\begin{equation}\label{ms:finite:five:eq:moments}
 m\ge\frac{1+\beta^2}{2\rho},\qquad
 \frac{2m^2}{1+\beta}\le p\le1-\frac{2m^2}{1-\beta},\qquad
 m^2\le pq.
\end{equation}
The first follows from $\rho\ge(a_+^2+a_-^2)/m$, the next two
from $m^2\le p a_+$ and $m^2\le q a_-$, and the last from
$1\ge m^2/p+m^2/q$.

Take a closed beta slab $[b_l,b_u]\subset[17/50,b_3]$ while retaining
the whole moment interval $\rho\in[477/400,6/5]$. Set
\[
 m_l=\frac{1+b_l^2}{2(6/5)},\qquad
 p_l=\frac{2m_l^2}{1+b_u},\qquad
 p_u=1-\frac{2m_l^2}{1-b_l}.
\]
Every actual law in the slab lies in this probability interval.
The scalar domain gates require $0<p_l\le p_u<1/2$ for every
nonempty slab. Round this interval outward only while preserving
$0<l\le u<1/2$. For a closed box $p\in[l,u]$, $z\in[a,b]$ with
$0\le a\le b\le5$, define
\begin{align*}
 \bar v_0&=\max\left\{0,u^2
             \left(\frac{(1-l)(1-b_l)}{2m_l^2}-1\right)\right\},\\
 \bar v_1&=\max\left\{0,(1-l)^2
             \left(\frac{u(1+b_u)}{2m_l^2}-1\right)\right\}.
\end{align*}
They bound \eqref{ms:finite:five:eq:v0}--\eqref{ms:finite:five:eq:v1} above: on an
actual law each bracket is nonnegative and every positive factor has
been enlarged. If an upper bracket is negative, the actual-law domain
is empty; clipping to zero remains a safe relaxation.
The maximum of $b_j(p)$ on $[l,u]$ is at
$p_j=\min(u,\max(l,j/5))$. Since $T$ increases with $V$ and
decreases with a positive distance, a whole-box upper CDF is
\[
 \bar Q=\min\left\{1,\sum_{j=0}^5b_j(p_j)
 T\bigl((5-j)\bar v_0+j\bar v_1,j-b\bigr)\right\}.
\]
This includes $j-b=0$ with the preceding closed-boundary convention.

Let $\nu=a-5u$. A lower bound for the Gaussian argument is
\begin{equation}\label{ms:finite:five:eq:normal}
 x_l=\begin{cases}
 \nu m_l/(\sqrt5\,u(1-u)),&\nu\ge0,\\
 \nu/\sqrt{5u(1-u)},&\nu<0.
 \end{cases}
\end{equation}
For the first branch, use the lower mean and the upper variance factor.
For the second branch, if $z-5p\ge0$ the negative lower bound is
automatic. Otherwise $m\le\sqrt{pq}$ reduces the claim to the
monotonicity, valid for $z\ge0$ and $p<1/2$,
\[
 \frac{d}{dp}\frac{z-5p}{\sqrt{p(1-p)}}
 =-\frac{5p+z(1-2p)}{2[p(1-p)]^{3/2}}<0.
\]
Monotonicity in $z$ then gives \eqref{ms:finite:five:eq:normal}.
For a negative numerator, exact evaluation divides by a lower square
root; for a nonnegative numerator it divides by an upper square root.
The resulting whole-box score bound is
\begin{equation}\label{ms:finite:five:eq:boxscore}
 \frac{\sqrt5}{477/400}\max\{0,\bar Q-\Phi(x_l)\}.
\end{equation}
The positive part must be taken before an upper-factor substitution;
otherwise a negative difference would reverse the required bound.

\subsubsection{Finite subdivision and the two infinite phase rays}

Use the closed beta slabs with successive endpoints
$0.340,0.341,\ldots,0.351,b_3$, subject to exact checking of their
order and the domain gates above. Probability intervals may be rounded
outward to a $10^{-6}$ grid when those gates remain true. The compact
phase cover is $[0,1],\ldots,[4,5]$. A certificate consists of finite
binary midpoint subdivisions in $p,z$, retaining both closed children.
At every leaf, rational Gaussian and radical enclosures must establish
that \eqref{ms:finite:five:eq:boxscore} is strictly below $c_0$.

For the exterior phase rays take the same moment bounds $m_L,p_L,p_U$
for the whole corner $\beta\in[17/50,b_3]$. With the same exact gates,
put
\[
 v_U=p_U(1-p_U),\qquad s_U=v_U^2/m_L^2-v_U,\qquad
 V_j^U=\max\left\{\frac j{p_L},\frac{5-j}{1-p_U}\right\}s_U.
\]
For an actual law $v=pq\ge m_L^2$ and
$s=\E W^2=v^2/m^2-v$. The function $v^2/m_L^2-v$ increases on
$v\ge m_L^2$, so $s\le s_U$. Also $qv_0+pv_1=s$ implies
$V_j\le\max\{j/p,(5-j)/q\}s\le V_j^U$, and
$b_j(p)\le\binom5j p_U^j(1-p_L)^{5-j}$.
For $z\le0$, drop the nonnegative Gaussian CDF and use monotonicity
at zero. The uniform score upper bound is
\begin{equation}\label{ms:finite:five:eq:lefttail}
 B_-:=\frac{\sqrt5}{477/400}
 \left[(1-p_L)^5+\sum_{j=1}^5\binom5j p_U^j(1-p_L)^{5-j}
             \frac{V_j^U}{V_j^U+j^2}\right].
\end{equation}
For $z\ge5$, use the upper CDF bound one. Its Gaussian argument is
at least its value at five, namely $\sqrt5\,m/p$, and hence at least
$\sqrt5\,m_L/p_U$. Cantelli for a standard normal gives
\begin{equation}\label{ms:finite:five:eq:righttail}
 B_+:=\frac{\sqrt5}{477/400}\frac1{1+5m_L^2/p_U^2}.
\end{equation}
The scalar certificate must prove $B_-<c_0$ and $B_+<c_0$ with
outward square roots. Thus no unbounded threshold range is discarded.

\begin{proposition}[Certificate-conditional gap-free corner]
\label{ms:finite:five:prop:gapfree}
Suppose the beta enclosures and domain gates above hold, a complete
finite compact-phase subdivision has every whole-box bound
\eqref{ms:finite:five:eq:boxscore} below $c_0$, and
\eqref{ms:finite:five:eq:lefttail}--\eqref{ms:finite:five:eq:righttail} are below $c_0$.
Then every standardized iid marginal with
\[
 \rho\in[477/400,6/5],\qquad\beta\in[17/50,b_3]
\]
has positive normalized closed-CDF error below $c_0$ at every raw
threshold. This implication requires no predecessor, support, contact,
or physical-gap hypothesis.
\end{proposition}
\begin{proof}
Each actual law belongs to a closed beta slab and its closed probability
interval by \eqref{ms:finite:five:eq:moments}. Each compact phase belongs to an
original box. Induction up its finite binary partition applies
\eqref{ms:finite:five:eq:boxscore} on every leaf, including shared cut faces.
The two exterior bounds cover the remaining phases. Since
$h>0$, the map $H\mapsto z=5p+hH$ covers the whole real line.
\end{proof}

\subsection{The acquired resource and count obligations at five}
\label{ms:finite:five:sec:acquired}

For the attained contrary law outside the direct corner, the implication
chain is contact and $q_5=0$, then support and Fourier bounds, then an
initial resource cover, then propagated resources, then counts.
The initial cover starts with no assumed physical defect or central-gap
seed. Every final surviving cell must acquire a genuine central support
gap before an inequality requiring that gap is applied.
The continuous statements required are the contact lemma and conserved identity
in Assumption~\ref{ms:finite:inputs}, items~\ref{ms:finite:input:cut49:global:contact}
and~\ref{ms:finite:input:afcore:prop:conserved}; the resource enclosure and
propagation in item~\ref{ms:finite:input:rescore:u-contract}; and the count
implication in Proposition~\ref{ms:finite:countcore:forest}. Their finite gates still have to be verified
on the actual fifth-index domains; an invocation by name is not such a
verification.

The same physical defect is used throughout:
\[
 D=(3+\sqrt{10})\rho-\gamma-3J,\qquad
 \gamma=\E X^3,\qquad J=\sum_{(a,b)}\tau_{a,b}(b-a)^2,
 \quad \tau(t)=\E[X\mathbf 1_{\{X\ge t\}}].
\]
The bounded complementary support intervals use the half-open convention
in Assumption~\ref{ms:finite:inputs}, item~\ref{ms:finite:input:afcore:eq:gap-definitions}. The conserved identity integrates
the same contact slack in two ways and retains the complete Fourier tail;
it does not use a large-index expansion. All formulas are evaluated at
$n=5$, with their positivity and denominator gates retained.

After a central gap $(-A,B)$ is acquired, the bounds concern the same
marginal's moments, endpoints and conditional noise. In particular,
already proved floors $m\ge m_L$, $A\ge a$, $B\ge b$ retain the price
$6mAB\ge6m_Lab$. It may be combined with another lower price by their
maximum, not by an unsupported sum. If $6mAB\ge P$ and $m\le m_U$,
then $4AB\le(A+B)^2$ gives $A+B\ge\sqrt{2P/(3m_U)}$.
For the new price alone,
\[
 \frac{2(6m_Lab)}{3m_U}\le4ab\le(a+b)^2.
\]
It therefore retains a proved gap and cannot create a gap from a failed
initial cover. The simultaneous signed-resource consequences are those
of Assumption~\ref{ms:finite:inputs}, item~\ref{ms:finite:input:rescore:u-three-uses}, with the same phase passed through
Assumption~\ref{ms:finite:inputs}, item~\ref{ms:finite:input:rescore:u-phase-interface}.

\subsubsection{The actual iid count probability}

Use the representation of Section~\ref{ms:finite:countcore:sec}. At nearest-count
phase $z=k+u$, let $W_j$ be the sum of the five residuals conditional on
exactly $j$ positive signs. The common closed-count identity gives
\begin{equation}\label{ms:finite:five:eq:outgoing}
 \Pp(S_5\le H)\le F_{\mathrm{Bin}(5,p)}(k)
 -b_k(p)\Pp(W_k>u)
 +\sum_{j>k}b_j(p)\Pp(W_j\le u-(j-k)).
\end{equation}
The outgoing event is strict and the incoming events are closed.
An atom on the event plane is never subtracted twice.
Incoming bounds use the proved Cantelli, exponential or stop-loss
inequalities of the common count section. A basic outgoing minorant is
\[
 G_a(x)=-x^4+3a^2x^2+2a^3x=x(x+a)^2(2a-x),\qquad a>0.
\]
It is nonpositive on $x\le0$ and has global maximum
$(9/4+3\sqrt3/2)a^4$. Thus
\[
 \Pp(U>u)\ge\frac{\E G_a(U-u)}{(9/4+3\sqrt3/2)a^4}.
\]
This is the $c=2$ case of Lemma~\ref{ms:finite:countcore:quartic}.
Acquired bounded support and common conditional moments control its
expectation. The same pair of conditional third moments is retained in
all constraints. Linear objectives, or affine ratios with a proved
strictly positive denominator, attain their maxima over the resulting
bounded polygon at its vertices. The polygon is a relaxation containing
every actual law; its feasible points need not represent convolution laws.

A finite count certificate must reconstruct each of the six original
roots $k=0,\ldots,5$, not just its contracted survivors. It must verify
all coupled moment, physical-defect, phase and Gaussian inequalities on
each original closed leaf box. A leaf closes only by a proved
necessary-condition contradiction or an upper score below a fixed
rational lower enclosure of $\CE$. The acquired mean, scale and
threshold bounds must justify the original nearest-count domain as in
Proposition~\ref{ms:finite:countcore:forest}; any thresholds outside it require
their own tail bounds. Zero noise and equality faces are included.

\subsubsection{The conditional fifth-index conclusion}

\begin{theorem}[Certificate-conditional fifth-index composition]
\label{ms:finite:five:thm:composition}
Assume $C_4\le\CE$. Suppose exact finite certificates establish all
of the following obligations using the continuous inequalities above:
\begin{enumerate}
\item the scalar beta and target enclosures, all domain gates, and the
      structural-tail endpoint comparison in Section~\ref{ms:finite:five:sec:cover};
\item quantitative support followed by the complete Fourier comparisons
      on every listed moment strip above $6/5$;
\item complete initial and propagated same-law resource covers on every
      resources/counts row, with either a proved contradiction or a
      complete six-root count certificate at each surviving panel;
\item the compact-phase forest and both infinite phase-ray comparisons
      in Proposition~\ref{ms:finite:five:prop:gapfree}.
\end{enumerate}
Then $C_5\le\CE$.
\end{theorem}
\begin{proof}
Suppose $C_5>\CE$. Lemma~\ref{ms:finite:five:lem:qzero} supplies an attained
positive closed-edge witness with zero predecessor allowance. Its
moments satisfy $\rho\ge1$ and \eqref{ms:finite:five:eq:beta}.
If $\rho\ge6/5$, the acquired support, complete Fourier cover and
decreasing structural tail exclude it by the first two certificate
obligations. Otherwise it belongs to a closed row of the low-moment
table. The third obligation and the continuous resource/count
implications exclude every resources/counts row. The fourth obligation
and Proposition~\ref{ms:finite:five:prop:gapfree} exclude the remaining corner.
All shared boundaries, in particular $\rho=477/400$, $\rho=6/5$
and $\beta=17/50$, remain included. Hence no contrary maximizer exists.
\end{proof}

The finite certificate conditions of
Theorem~\ref{ms:finite:five:thm:composition} remain to be completed.

\subsection{Signed moment and physical-gap refinements}
\label{ms:finite:signed-refinements}

\subsubsection{A sign-aware paired Fourier coefficient rectangle}
\label{ms:finite:paired-rectangle}
Assume $\E X=0$, $\E X^2=1$, $|X|\le R$, and
\[
 \mu_4\le M\le111/10,\quad |\gamma|\le g\le37/20,
 \quad \E|X|^5\le\eta_5,\quad
 0\le g_l\le|\gamma|,\quad M\ge1+g_l^2.
\]
Here $R\ge1$. Reflection is permitted only inside the auxiliary Fourier
integral, to write $t=|\gamma|\ge0$; the signed $\gamma$ in the actual
contact identity is retained. Put $Q=X^2-tX-1$ after this auxiliary
reflection. The Pearson identities
\eqref{ms:fourier:rescore:fourier-pearson}, applied to this same
standardized law, give
\begin{align}
 \E Q^2&=\mu_4-1-t^2=:v_Q\ge0,\notag\\
 \mu_5&=t^3+2t+\E[Q^2(X+2t)],\notag\\
 \mu_6&=1+3t^2+t^4+
   \E[Q^2((X+t)^2+2t^2+2)].
 \label{ms:finite:pearson-identities}
\end{align}
For $z=(10t-\mu_5)/120$ and $c=(3-\mu_4)t/144$ this proves
\begin{equation}
 \begin{gathered}
 L(t)\le120z\le U(t),\\
 L(t)=8t-t^3-(R+2t)(M-1-t^2),\\
 U(t)=8t-t^3+(R-2t)_+(M-1-t^2).
 \end{gathered}
 \label{ms:finite:paired-z}
\end{equation}
Indeed $-R+2t\le X+2t\le R+2t$. For the upper bound on $z$, the
coefficient $R-2t$ multiplies $v_Q$: use $v_Q\le M-1-t^2$ if this
coefficient is nonnegative, and $v_Q\ge0$ otherwise. The lower bound
uses the nonnegative coefficient $R+2t$. This explains the sign switch
without requiring $R>2g$.

Every coefficient enclosure must cover the whole interval $[g_l,g]$.
The cubic
\[
 L(t)=-R(M-1)+(10-2M)t+Rt^2+t^3
\]
has its minimum at an endpoint or a root of
$3t^2+2Rt+10-2M$ in the interval; if $M\le5$ it is increasing.
For $U$, split at $R/2$ and maximize the cubics
\[
 R(M-1)+(10-2M)t-Rt^2+t^3,\qquad 8t-t^3
\]
on their respective closed subintervals. Endpoints and derivative roots
suffice. Retaining infeasible $t$ with $t^2>M-1$ merely enlarges the
enclosure domain; \eqref{ms:finite:paired-z} is used only for actual laws.
If $\ell\le\min L$ and $u\ge\max U$, safe endpoints are
\[
 z_l=\max\{-\eta_5-10g,\ell\}/120,\quad
 z_u=\min\{\eta_5+10g,u\}/120,\quad
 Z=\max(|z_l|,|z_u|),
\]
with $z_l\le z_u$ required. The deliberately symmetric absolute fifth
moment bound in this intersection is valid even though $t\ge0$.
The other coefficients satisfy
\[
 |(3-\mu_4)/24|\le a_*:=\max\{|3-M|,|2-g_l^2|\}/24,
\]
\[
 c_l=\frac{3-M}{144}
 \begin{cases}g_l,&M\le3,\\g,&M>3,\end{cases}
 \qquad c_u=\frac1{144}\max_{g_l\le t\le g}(2t-t^3).
\]
These follow from $1+t^2\le\mu_4\le M$. The maximum in $c_u$ is
at $g$, at $g_l$, or at $\sqrt{2/3}$ according as the interval is
below, above, or contains that stationary point.

For completeness, the cancellation norm needed by the paired transfer is
\[
 J(z,c)=\int_0^\infty e^{-x^2/2}|zx^2-cx^4|\,dx.
\]
The cancellation-norm calculation in the proof of
Lemma~\ref{ms:fourier:rescore:fourier-early-low} applies to every real
pair $(z,c)$: $J$ is convex, its enclosing rectangle is bounded by its
four corners, and $J(-z,-c)=J(z,c)$. When $zc\le0$,
$J=\sqrt{\pi/2}(|z|+3|c|)$. When $zc>0$, reverse both signs
if necessary and put $a=\sqrt{z/c}$. The exact Gaussian-tail identity
in the proof of that lemma becomes
\begin{equation}
 J(z,c)=\sqrt{\pi/2}(z-3c)(1-4\{1-\Phi(a)\})
          +6ca e^{-a^2/2},\qquad z,c>0.
 \label{ms:finite:paired-mills}
\end{equation}
Here $\int_a^\infty e^{-x^2/2}\,dx=\sqrt{2\pi}(1-\Phi(a))$,
so this formula retains the signed tail term before any upper bound.
In particular a possibly negative coefficient $z-3c$ requires the
appropriate lower or upper Gaussian enclosure; taking its absolute value
is a different bound.

The sixth moment follows directly from \eqref{ms:finite:pearson-identities}:
\[
 \mu_6\le M_6:=\min\{R\eta_5,
 1+3g^2+g^4+(R^2+2Rg+3g^2+2)(M-1-g_l^2)\}.
\]
All multiplied factors are nonnegative. These coefficient bounds may be
inserted into a separately proved signed Fourier transfer only on its
original domain ($n>2$ and positive damping included). Its damping is
$b_0=(45-M)(1-2/n)/96$, not $43(1-2/n)/96$.
They establish no Fourier integral or cutoff by themselves.

\begin{remark}[Why the old sign substitution cannot be widened]
For $\Pp(X=3/2)=4/13$ and $\Pp(X=-2/3)=9/13$,
$\gamma=5/6$, $\mu_4=61/36$ and $120z=1315/216$.
The valid loose bounds $R=3/2$, $g=7/5$, $M=197/40$ obey $R>g$,
but the old unsigned quadratic value $951/160$ is smaller than
$1315/216$. Thus dropping $R\ge2g$ from that older formula is invalid;
the piecewise argument above is a distinct repaired coefficient bound.
\end{remark}

\subsubsection{A quadratic signed-moment loss}
\begin{lemma}\label{ms:finite:signed-loss}
For a standardized law with finite fifth absolute moment, write
$\rho=\E|X|^3$, $\beta=\E X|X|$, $\gamma=\E X^3$ and
$m=\E X_+>0$. Then
\[
 \frac{\rho\E X^4}{2}-\frac{\E|X|^5}{4}
 \le\frac{\rho^3}{4}-\frac{(\gamma-\rho\beta)^2}{8m}.
\]
Equality holds for every centered standardized two-point law.
\end{lemma}
\begin{proof}
Use the weighted Cauchy--Schwarz calculation leading to
\eqref{ms:fourier:rescore:eq:weighted-constant-loss}, with
$Y=|X|$ and $Z=Y^2-\rho Y$. That calculation requires only
$\E|X|^5<\infty$ and $\E|X|=2m>0$, which hold here; it also
includes a possible zero atom.
For support $\{-a,b\}$ with $ab=1$,
$\rho=(a^2+b^2)/(a+b)$ and $a^2-\rho a=-(b^2-\rho b)$.
Thus $Z$ is proportional to $\operatorname{sgn}(X)$ on the support,
which is precisely the equality condition.
\end{proof}
On a rectangular $(\rho,\beta,\gamma)$ domain, all eight corners
enclose $\delta=\gamma-\rho\beta$. If its interval has distance
$d_\delta$ from zero and $m\le m_u$, the certified loss is
$d_\delta^2/(8m_u)$. It is subtracted only in this matching constant
term. It does not by itself bound the physical gap defect.

\subsubsection{Sign probabilities, including zero atoms}
\begin{lemma}\label{ms:finite:sign-probability}
Every standardized law with finite $\rho$ satisfies
\[
 \Pp(X>0),\ \Pp(X<0)\ \ge\ \kappa/\rho^2,
 \qquad \kappa=\frac8{27}(7\sqrt7-17).
\]
Consequently $p=\Pp(X>0)\in[\kappa/\rho^2,1-\kappa/\rho^2]$
even when $\Pp(X=0)>0$.
\end{lemma}
\begin{proof}
Put $a=\E X_+^2\in(0,1)$ and $m=\E X_+=\E X_->0$.
Signwise Cauchy gives
$a^2\le m\E X_+^3$, $(1-a)^2\le m\E X_-^3$ and $m^2\le pa$.
Therefore $p\rho^2\ge(2a^2-2a+1)^2/a$. Its derivative is
\[
 \frac{(2a^2-2a+1)(6a^2-2a-1)}{a^2}.
\]
The unique minimum in $(0,1)$ is $a_0=(1+\sqrt7)/6$, where the
value is $\kappa$. Reflect to obtain the negative-sign inequality;
$\Pp(X<0)\le1-p$ proves the upper bound on $p$.
Sharpness follows by fixing $a=a_0$, $0<p<1-a$, $m=\sqrt{pa}$,
and assigning masses $p$ at $a/m$, $pa/(1-a)$ at $-(1-a)/m$,
and the remaining mass at zero. Direct substitution gives mean zero,
variance one, and $p\rho^2=\kappa$.
\end{proof}
The rational lower constant $11261/25000$ is licensed by
$\sqrt7>10583/4000$ (square both positive sides), followed by
rational substitution. No support or optimizer hypothesis enters this lemma.

\subsubsection{Spending one outer-gap residual}
For a standardized law $P$, assume a genuine central support gap
$(-A,B)$ with $A,B>0$, and $\supp P\subset[-\mathcal R_-,\mathcal R_+]$. Use the same physical
$D$ as in Lemma~\ref{ms:physical-defect}, and write
$m=\E X_+=\E X_-$, $\sigma=\sqrt{10}$. The moment residuals in that
lemma are
\[
 T_+=\E X_+^3-mB^2,\quad R_+=\E[X_+(X_+-B)],
 \qquad T_-=\E X_-^3-mA^2,\quad R_-=\E[X_-(X_--A)].
\]
On the positive support,
$X(X^2-B^2)\le(\mathcal R_++B)X(X-B)$, hence
$T_+\le(\mathcal R_++B)R_+$; reflection gives the other inequality.
All four residuals are nonnegative. Lemma~\ref{ms:physical-defect}
and \eqref{ms:defect-decomposition} give
\begin{align}
 D&=mQ+(\sigma-1)T_++(\sigma+1)T_-+3(T_++T_--J_o),\notag\\
 Q&=(\sqrt{\sigma+1}A-\sqrt{\sigma-1}B)^2,
 \qquad T_++T_--J_o\ge2BR_++2AR_-.
 \label{ms:finite:outer-decomposition}
\end{align}
For acquired lower endpoints $A_l,B_l>0$ and $S_l\le\sqrt{10}$,
\begin{equation}
 w_+T_++w_-T_-\le D-mQ,\quad
 w_+=S_l-1+\frac{6B_l}{\mathcal R_++B_l},\quad
 w_-=S_l+1+\frac{6A_l}{\mathcal R_-+A_l}.
 \label{ms:finite:outer-weights}
\end{equation}
Require $w_\pm>0$. The ratios increase in the endpoint and decrease
in the support radius, proving the replacement directions.
If $Q\ge q_m\ge0$, $m\ge m_l$, $D\le d$, the common residual is
$d_o=d-m_lq_m$. A negative $d_o$ excludes the cell. Otherwise
\[
 T_+\le\min\{B_{3,+}-m_lB_l^2,d_o/w_+\},\qquad
 T_-\le\min\{B_{3,-}-m_lA_l^2,d_o/w_-\},
\]
where $B_{3,\pm}$ are separately acquired third-moment upper bounds.
If the retained integral bound is
$\mathrm{outer}\le(3/8)(k_+T_++k_-T_-)$, then
\[
 \mathrm{outer}\le\frac38\max(k_+/w_+,k_-/w_-)(D-mQ).
\]
In the application below $k_+=(\rho_u-B_l)_+^2$ and
$k_-=(\mathcal R_--\rho_l)_+^2$. The implication follows term by
term from \eqref{ms:finite:outer-weights}; it does not spend two
independent copies of $D-mQ$. The integral bound itself remains a
retained analytic input when this lemma is applied.

\subsubsection{Gap-localizing inequalities and mean contraction}
\begin{lemma}\label{ms:finite:gap-localizing}
Suppose the positive support is at least a proved $B_l>0$. Set
$A_2=\E X_+^2$, $B_3=\E X_+^3$, $m=\E X_+$ and $p=\Pp(X>0)$.
Then
\[
 (m-B_lp)(B_3-B_lA_2)\ge(A_2-B_lm)^2.
\]
If $m\le m_u$, $A_2\ge a_{2,l}$ and $B_3\le b_{3,u}$, put
$d=b_{3,u}-B_la_{2,l}$. For $d>0$,
\[
 p\le\frac{m_u-(a_{2,l}-B_lm_u)_+^2/d}{B_l}.
\]
For $d<0$ the cell is empty. For $d=0$ a positive lower numerator
excludes it, and otherwise $p\le m_u/B_l$ is valid.
\end{lemma}
\begin{proof}
Apply Cauchy--Schwarz to $1$ and $X$ under the nonnegative measure
$(X-B_l)\ind_{X>0}\,dP$. All three entries in its moment matrix are
nonnegative. Enlarge the second diagonal entry to $d$ and lower the
off-diagonal entry to $(a_{2,l}-B_lm_u)_+$. Rearrangement gives the
positive-denominator case; a nonnegative diagonal cannot be below zero,
and a zero diagonal forces the off-diagonal entry to vanish, proving the
other cases.
\end{proof}
For fixed endpoints and moment bounds, the resulting cap $P_+(m)$ is
nondecreasing, since where $d>0$ its derivative is
$[1+2B_l(a_{2,l}-B_lm)_+/d]/B_l>0$. The reflected cap $P_-(m)$
has the same property. Thus the necessary conditions
$P_+(m)\ge p_l$, $P_-(m)\ge1-p_u$, and $P_+(m)+P_-(m)\ge1$
are monotone. Failure at $m_u$ excludes the cell; finite bisection
retaining an infeasible lower endpoint produces a valid lower bound on
every feasible $m$, without an assumption of convergence. None of these
arguments acquires the central gap it assumes.

\subsubsection{Exact Gaussian phase with an affine physical-defect price}
\label{ms:finite:exact-phase}
Consider $n\in\{6,7,8,9\}$. Fix one positive-branch
maximizing law, a raw threshold with a previously proved
$|H|\le H_{\max}$, and its score $C$, moments $\rho,\gamma$, and
physical defect $D\ge0$. Assume the complete conserved inequality
\begin{align}
 \frac{n\phi D}{6}+n(C-\CE e^{-H^2/(2n)})\rho
 \le K_F+s_F|H|+C A_{\rm act}
 +H(CK+\phi U)-H^2(C\rho+\gamma\phi/2),
 \label{ms:finite:phase-input}
\end{align}
where $\phi=\phi(H/\sqrt n)$, $K_F,s_F\ge0$ include complete
Fourier costs, and every omitted slack is nonnegative. Assume acquired
intervals for $C>0,\rho\ge1,K,U$, a lower $\gamma_l$, and
$A_{\rm act}\le A_*-t_DD$ with $t_D\ge0$ independent of $C$.
For $H=\epsilon u$, $u\ge0$, choose $(K_\epsilon,U_\epsilon)$ to
be the upper endpoints if $\epsilon=1$ and the lower endpoints if
$\epsilon=-1$. Substitution and $\CE/\phi(0)=(3+\sqrt{10})/6$ give
\begin{equation}
 D\le\frac{P(u)+\phi Q(u)}{n\phi/6+Ct_D},
 \label{ms:finite:phase-ratio}
\end{equation}
where
\begin{align*}
 P(u)&=K_F+CA_*-nC\rho+(s_F+\epsilon CK_\epsilon)u-C\rho u^2,\\
 Q(u)&=n(3+\sqrt{10})\rho/6+\epsilon U_\epsilon u-\gamma_lu^2/2.
\end{align*}
The denominator is strictly positive. For fixed $u,\phi$, the ratio
is affine in $\rho$; for fixed $\rho$ it is fractional-affine in
$C$ with positive denominator. Its derivative in $C$ has constant
sign. Thus all four $(C,\rho)$ corners suffice, without asserting that
the corners themselves are realizable laws.

Cover $[0,L]$, $L=\lceil MH_{\max}\rceil/M$, by all closed mesh
cells of width $1/M$ and evaluate both signs. On each cell the exact
maxima $P_*,Q_*$ occur at endpoints or an interior stationary point
of a concave quadratic. With $\phi$ enclosed by its positive endpoint
density bounds, replacing $P,Q$ by $P_*,Q_*$ enlarges the ratio.
It is again fractional-affine in $\phi$, so its two endpoints suffice.
The maximum over the complete mesh bounds $D$ above; a negative
maximum is a contradiction. Application requires the conserved input \eqref{ms:finite:phase-input}
and the complete threshold bound $H_{\max}$ at the stated actual index.
The finite Gaussian enclosure additionally requires $L^2/(2n)\le7$.

\subsubsection{The central signed price may be negative}
Suppose an already acquired genuine central gap gives, for the same law,
\[
 \begin{gathered}
 A_{\rm act}\le\tfrac32\beta(J+\gamma)+\mathrm{central}
 +\mathrm{outer}+\rho^3/4-\rho/2-\mathrm{penalty},\\
 \mathrm{outer}\le\lambda_o(D-mQ),\qquad mQ\ge q_l\ge0.
 \end{gathered}
\]
Here $\lambda_o\ge0$, the central term is an upper bound, and the
penalty is a lower bound. Substitute
$J=((3+\sqrt{10})\rho-\gamma-D)/3$ to obtain
\begin{align}
 A_{\rm act}&\le A_{\rm cen}-t_{\rm cen}D,\notag\\
 t_{\rm cen}&=\beta/2-\lambda_o,\notag\\
 A_{\rm cen}&=\beta((3+\sqrt{10})\rho/2+\gamma)
  +\mathrm{central}-\lambda_oq_l+\rho^3/4-\rho/2-\mathrm{penalty}.
 \label{ms:finite:central-price}
\end{align}
The first term is multiaffine in its rectangular variables and is
bounded at all corners; use $\rho_u^3/4-\rho_l/2$ for the remaining
moment term. A lower price is $\beta_l/2-\lambda_o$, since $D\ge0$.
This price can be negative. In that case the proof of
\eqref{ms:finite:phase-ratio} still works only if its \emph{total}
denominator is proved positive on the complete domain. Different
complete candidates may be minimized after each has maximized over
its entire score interval. They may not spend the common defect twice.


\subsection{Symmetrization and sine inputs from one marginal}
\label{ms:finite:symmetrization}
Here $\eta_s,\zeta_s$ are symmetrization parameters, distinct from the
four-summand excess $\eta$. In this section $h=\E|X|$, distinct from
the count scale previously called $h$.

\subsubsection{Fixed-first-moment envelopes}\label{ms:finite:fixed-h}
Let $0<p\le1/2$, $q=1-p$, $v=pq$, and write the same marginal as
$Y=-p+U_0$ on mass $q$ and $Y=q+U_1$ on mass $p$, with $\E U_i=0$.
Put $b=q\E U_0^2$, $c=p\E U_1^2$, $s=b+c$, $w=v+s$,
$X=Y/\sqrt w$, and $h=2v/\sqrt w$. Assume acquired inward bounds
$U_0\le a_0$, $-U_1\le a_1$, $a_0+a_1\le d\le p$. Define, for
an independent copy with the same marginal,
\[
 \eta_s=\tfrac12\E|X-X'|^3-\rho,\qquad
 \zeta_s=4\E X_+^2\E X_-^2/h.
\]
The scalar identity
\[
 |x-y|^3=(x-y)^2(|x|+|y|)
 -2\min(|x|,|y|)(x-y)^2\ind_{xy>0}
\]
gives $\eta_s=h-\E[\min(|X|,|X'|)(X-X')^2\ind_{XX'>0}]$:
the expectation of the first term is $2\rho+2h$ by independence and
centering. On a same-sign pair, the smaller magnitude is at least
$(p-a_0)/\sqrt w$ or $(q-a_1)/\sqrt w$, while the conditional
expected squared difference is $2\E U_i^2/w$. The deficit is at least
\[
 \frac{2\{qb(p-a_0)+pc(q-a_1)\}}{w^{3/2}}
 \ge\frac{2s(v-qd)}{w^{3/2}},
\]
since $qa_0b+pa_1c\le q(a_0+a_1)s$. Consequently
\begin{equation}\label{ms:finite:eta-fixedh}
 \eta_s\le\frac{2(v^2+sqd)}{(v+s)^{3/2}}
 =E(p,h,d):=\frac{h^3}{4v}(1-d/p)+(d/p)h.
\end{equation}
This includes $s=0$ and sign-boundary atoms. The equality follows by
$s=4v^2/h^2-v$. For fixed $p,d$, $E$ increases with $h>0$. For fixed
$h$, put $a=h^2/(4v)$. On $a\le1$,
\[
 E_p/h=-a(1-2p)(1-d/p)/v-d(1-a)/p^2\le0,
 \qquad E_d=h(1-a)/p\ge0.
\]
Thus $p_l\le p\le p_u$, $h\le H$, $H^2\le4p_l(1-p_l)$ and
$d\le d_u\le p_l$ imply $\eta_s\le E(p_l,H,d_u)$. If the squared
gate fails, retain $\eta_s\le h\le H$; the derivative argument is
not used at virtual points with $a>1$.
For $0<H\le1$ and $v_l=p_l(1-p_l)\ge1/5$, $4v^2/H^2-v$ increases
on $v\ge v_l$. Hence $s\ge s_L=\max(0,4v_l^2/H^2-v_l)$.
If $s_L>s_u$ the cell is empty; otherwise an acquired budget
$d(s)\le\min(A-Bs,C-Ds)$, $B,D\ge0$, supplies
$d_u=\min(A-Bs_L,C-Ds_L)$. A negative $d_u$ excludes the cell;
it is not silently changed to zero.

The signwise second moments are $(vq+c)/w$ and $(vp+b)/w$, so
$\zeta_s=2(vp+b)(vq+c)/(vw^{3/2})$. Set $b=s(1-\alpha)$,
$c=s\alpha$. The numerator is concave in $\alpha$; its smaller
endpoint for $p\le q$ is $\alpha=1$. Therefore
\begin{equation}\label{ms:finite:zeta-fixedh}
 \zeta_s\ge\frac{2(v^2+ps)}{(v+s)^{3/2}}
 =Z(p,h):=h/q-h^3/(4q^2),\qquad \zeta_s\ge h.
\end{equation}
For the second bound use
$(vp+b)(vq+c)-v^2(v+s)=v(p^2c+q^2b)+bc\ge0$.
For $p\le43/100$ and $0<h\le1$, $Z_p=h(q-h^2/2)/q^3>0$ and
$Z_{hh}=-3h/(2q^2)<0$. Thus $L\le h\le H\le1$ yields
$\zeta_s\ge\max\{L,\min(Z(p_l,L),Z(p_l,H))\}$ by concavity on
the entire interval. If $L^2\ge4(1-p_l)/3$, $Z(p_l,h)$ decreases
throughout and the minimum is $Z(p_l,H)$. These are moment inequalities;
the sine-circle theorem that consumes $\zeta_s$ is a separate input.

\subsubsection{The Bernoulli-shaped envelope with all its gates}
\begin{proposition}\label{ms:finite:bernoulli-envelope}
Assume the same sign coordinates, $8/25\le p\le43/100$,
$0\le s\le13/1000$, and $11/10\le\rho\le6/5$. Let fixed
$0<D\le9/100$, $2\le C\le3$, $Q\ge2/5$ satisfy
\[
 a_0+a_1\le D,\quad
 \rho\le U(v,s):=\frac{v(1-2v)+Cs-s^2/Q}{(v+s)^{3/2}},
 \quad C+75D/16<3.
\]
For an acquired $p_l\le p$ require
$8p_l-5p_l^2-2-(1-p_l)D\ge0$. Then
\[
 \eta_s\le B(\rho)\le\zeta_s,\qquad
 B(r)=(\sqrt{r^2+8}-r)/2.
\]
\end{proposition}
\begin{proof}
The centered one-sided cubic inequality gives
\[
 \rho\ge I(p,s,D):=
 \frac{v(1-2v)+(3p-D)s+s^2/(qD)}{(v+s)^{3/2}}.
\]
Indeed the inward cubic estimates give quadratic terms
$b^2/(qa_0)+c^2/(pa_1)\ge s^2/(qa_0+pa_1)\ge s^2/(qD)$
by Cauchy--Schwarz; their linear terms are at least $(3p-D)s$.
A zero inward cap forces the corresponding conditional variance to vanish,
so that term is omitted. The outward upper $U$ is a separate acquired
premise, not a consequence of the parameter rectangle.
Put $w=v+s$, $M_1=v^2+ps$, and let $N_1$ be the numerator of $I$.
With $Z=2M_1/w^{3/2}$, expansion gives
\begin{align*}
 (IZ+Z^2-2)w^3/2
 &=M_1(N_1+2M_1)-w^3\\
 &=svp[8p-5p^2-2-qD]+s^2p[pq/D+8p-3-D]
   +s^3[p/(qD)-1].
\end{align*}
The first bracket increases with $p$ (derivative $8-10p+D>0$) and
is nonnegative by the gate. The second is at least
$(1/5)/(9/100)+8(8/25)-3-9/100=1523/900>0$; the last is at least
$647/153>0$. Since $\rho\ge I$ and $\zeta_s\ge Z>0$, the positive
root of $x^2+\rho x-2=0$ proves the $\zeta_s$ assertion.

For the other assertion put $a=Dq$ and $E=2(v^2+as)/(v+s)^{3/2}$.
Up to positive denominators the derivatives of $U$ have numerators
\begin{align*}
 J_s&=v(C-3/2+3v)-(C/2+2v/Q)s-s^2/(2Q),\\
 J_v&=-v/2-v^2+s(1-4v-3C/2)+3s^2/(2Q).
\end{align*}
On the larger rectangle $1/5\le v\le1/4$, $0\le s\le1/40$,
$J_s\ge963/6400>0$ and $J_v\le-881/6400<0$. For $J_s$, bound
its first term below by $11/50$ and its negative terms in magnitude
by $11/160$ and $1/1280$. For $J_v$, drop its negative linear term
and use $-1/10-1/25+3/1280=-881/6400$.
At fixed $a$ the derivative numerators of $E/2$ are
$K_s=av-3v^2/2-as/2<0$ and $K_v=v^2/2+(2v-3a/2)s$,
since $a/v=D/p\le9/32<3/10$. Expansion gives
$K_vJ_s-K_sJ_v=\sum_{i=0}^3T_i s^i$, where
\begin{align*}
 T_0&=(v^2/2)[v(C-3)+a(1+2v)],\\
 T_1&=-(C/2+3/2)v^2+av(1-v)-v^3/Q,\\
 T_2&=-Cv+a(1/2-2v)+(-2v^2+3av/2)/Q,\\
 T_3&=(3a/2-v)/Q.
\end{align*}
All are negative. For $T_0$ use
$(a/v)(1+2v)\le75D/16<3-C$; the others use $a/v<3/10$,
$C\ge2$, $v\ge1/5$ and $1/2-2v\le1/10$.
Since $v'=1-2p$, $a'=-D$, the determinant for the actual $p$
derivative is
$(1-2p)\sum T_i s^i-Ds(v+s)J_s<0$. Hence $E$ decreases in
$p$ along a level $U=R$, and decreases in $s$ at fixed $p$.

To justify the complete level path, fix $R=\rho\in[11/10,6/5]$.
There is a unique $p_B\in(0,1/2)$ with
$(1-2v_B)/\sqrt{v_B}=R$, $v_B=p_B(1-p_B)$, and
$2\sqrt{v_B}=B(R)$. The left side decreases; its square at $8/25$
is $124609/85000>(6/5)^2$, so $p_B>8/25$.
If actual $p\le p_B$, $\eta_s\le h\le2\sqrt v\le B(R)$.
Otherwise, for every $v'\in[v_B,v]$, monotonicity of $U$ gives
$U(v',0)\le R\le U(v',s)$. There is a unique root
$s_R(v')\in[0,s]$; the intermediate value theorem and $J_s>0$
give a differentiable path entirely inside the checked rectangle. Holding
$D,C,Q$ fixed on that virtual path,
\[
 \eta_s\le E(p,s)\le E(p,s_R(v))\le E(p_B,0)=B(R).
\]
No support acquisition at a virtual $p$ is being asserted.
\end{proof}
Since $B$ decreases and $r+B(r)$ increases (differentiate their formulas),
on $\rho\in[R_l,R_u]$ the valid conclusions are
$\eta_s\le B(R_l)$, $\zeta_s\ge B(R_u)$ and
$\rho+\eta_s\le R_u+B(R_u)$.

\subsubsection{Literal finite-frequency and strip interfaces}
\label{ms:finite:sine-interface}
The comparison is specified at $n=26,\ldots,30$ on eleven probability
panels (indices $32,\ldots,42$), hence 55 acquired B1 rows. From each
same-law row it recomputes, with positive divisor $g$,
\[
 s_u=\min(\lambda_{\max},\lambda_{\rm acquired})/n,
 \quad D=\min(d_0,(d_{\rm budget}-b)/g),
\]
\[
 C=\max\{3p_u+O_0,3(1-p_l)+O_1\},\quad
 Q=\max_{p\in\{p_l,p_u\}}\{(1-p)O_0+pO_1\},
\]
and requires every gate of Proposition~\ref{ms:finite:bernoulli-envelope}.
The rows' acquisition and these numerical gates remain conditions.

The additional support step uses $a_0+a_1\le D\le q-p$ and $D<p$:
\[
 (p-a_0)^2+(q-a_1)^2\ge p^2+(q-D)^2.
\]
For a fixed used inward sum, moving inward mass from the second endpoint
to the first has derivative $2(q-p-D)+4a_0\ge0$; the minimum is
at $a_0=0$, and using all of $D$ decreases the squared sum.
The right side decreases in $p$. Thus a sufficient standardized lower
bound is $[p_u^2+(1-p_u-D)^2]/[p_u(1-p_u)+s_u]$.
Each row must prove this is at least $17/10$, the squared sum of its
outer radii is at most $1037/100$, and each radius is at most $13/5$.

Let 23 consecutive closed $\rho$ strips cover $[11/10,6/5]$. At each
upper endpoint $r_u$, rational $0<z_l\le z_u$ satisfying
$z_l^2+r_uz_l\le2\le z_u^2+r_uz_u$ enclose $B(r_u)$.
The simultaneous valid inputs are
\[
 \rho+\eta_s\le r_u+z_u,\qquad
 \rho^2-\zeta_s^2\le r_u^2-z_l^2.
\]
The number $z_u$ alone is not a uniform upper bound on $\eta_s$ over
the strip. The separate sine theorem may supply the improved factor
\[
 1-\frac{17u^2}{200}\left(1-\frac{1037u^2}{4200}\right)
\]
only under $13u/5\le\pi$ and $1037u^2/100\le42$, retaining the
old valid circle outside. Moment and rectangle-convexity gates remain.
\begin{condition}[B1 cover at the literal five indices]
\label{ms:finite:b1-cover}
All 55 same-law resource rows satisfy the preceding gates. At each of the
five actual indices and each of the 23 strips, a complete four-term Fourier
comparison with retained compatible profile $62/625$ proves score strictly
below $4097/10000$. All low/high mesh cells, endpoint pieces, infinite
normal tail, analytic hypotheses and strict comparisons are verified.
\end{condition}
There are 115 integral obligations. Under this condition every actual row
and every moment in it belongs to a closed strip and is excluded by its
whole-strip comparison, provided resource acquisition and the applicable
sharp predecessors hold. Every unresolved cell rejects. The conclusion is
restricted to B1 at $n=26,\ldots,30$.
\subsection{Conditional reuse of finite floor-index resource sequences}
\label{ms:finite:floor-transfer}
Fix a complete finite sequence of upper bounds evaluated at an index
floor $N$. The actual index is denoted by $m$; the first positive mean is denoted by $\mu=\E X_+$. The same
unrestricted maximizing law, its raw threshold $H$ at a positive closed-CDF
maximum, signed moments, support and conditional noises must be used throughout.
The maximizing laws at different indices may differ.

\subsubsection{The weighted invariant and a finite induction}
\label{ms:finite:weighted-induction}
Let $D=(3+\sqrt{10})\rho-\gamma-3J\ge0$ be the physical defect,
$\mathcal R$ the nonnegative complete raw Fourier remainder, and let
$r,d$ be saved upper entries. The required induction invariant is
\begin{equation}\label{ms:finite:weighted-invariant}
 \mathcal R_{\rm actual}\le r,\qquad mD_{\rm actual}\le Nd.
\end{equation}
Thus $D_{\rm actual}\le d$ when $m\ge N$, but the converse does
not license a weighted noise allowance. A charged phase budget $B$ with
$mD\le B/\kappa$, $\kappa=997/15000$, exports
$d=B/(\kappa N)$. A negative $B$ excludes the law before division.
Taking the minimum with a previous valid weighted entry, upward rounding,
and taking a maximum across a complete cover all preserve
\eqref{ms:finite:weighted-invariant}. Index-free moment, mean, sign,
gap and endpoint contractions may use the weaker $D\le d$ and still
refer to the same law; they do not optimize again within a parameter cell.
Induction over a finite saved list therefore justifies every successive
entry once each transition has the properties proved below. No convergence
or limiting argument is used.

Write $v=p(1-p)$, $s=v\theta/(1-\theta)$, and let the actual noise be
$\lambda_m=ms$. If a floor affine resource is interpreted as
\begin{equation}\label{ms:finite:scaled-charge}
 D\ge b+g d_{\rm in}+(N/m)h_N\lambda_m,
 \qquad b,g,h_N\ge0,
\end{equation}
then multiplying by $m$ and using $mD\le Nd$ and $m\ge N$ gives
$b+g d_{\rm in}+h_N\lambda_m\le d$.
This proves the original joint intercepts for the actual noise; it does
not identify $\lambda_m$ with the floor noise. Equivalently a new count
interface uses actual $n=m$, defect bound $(N/m)d$, charge $(N/m)h_N$,
and phase normalization $15000/(997m)$. Merely replacing the old integer
index is insufficient. A saved harmonic noise-geometry bound may instead
be reused through
\[
 \theta_m\le\frac1{1+mv_{\min}/\Lambda+v_{\min}/(B d_0)}
 \le\theta_N,
\]
with its positive geometry constant $B$; it is not a bound on
$\lambda_m$ itself. On an actual noise cell $[a,b]$, the simultaneous
consequences $\theta\le b/(Nv_{\min}+b)$ and
$d_{\rm in}\le d_0(1-a/\Lambda)$ retain the same resource budget.

\subsubsection{Index-free refinements and decreasing analytic envelopes}
\label{ms:finite:floor-envelopes}
The signwise Cauchy bounds
\[
 \mu^2\le p_u(1+\beta_u)/2,\quad
 \mu^2\le(1-p_l)(1-\beta_l)/2,\quad
 \mu^2\le\max_{p\in[p_l,p_u]}p(1-p)
\]
are independent of $m$. So are moment feasibility, acquired central-gap
bounds and physical endpoint prices. In particular, once separately proved
endpoint inequalities supply
\[
 P=\max\{(S_l-1)\rho_l,S_l\rho_l-\gamma_u\}-d>0
 \quad\hbox{and}\quad 6\mu AB\ge P,
\]
then $4AB\le(A+B)^2$ and $\mu\le\mu_u$ give
$(A+B)^2\ge2P/(3\mu_u)$. The older independent lower bound
$2\rho_l-2d/(S_l-1)$ may be combined by taking their maximum.
An INITIAL mean-distance bound $(\rho_l-d/(S_l-1))/\mu_u$ likewise
has no sample-index dependence. These comparisons require a genuine
gap, positive mean, and the high-frequency distance condition $d_0\le12$;
they cannot create a gap from a failed unseeded initial cover.

The Cantelli event in the high-frequency bound uses the same-law fourth-moment cap $M$ and a fixed
$\delta$, with $V=2M+2$, $h_e=2-\delta^2>0$,
$a=h_e^2/(V+h_e^2)$ and $m_e=2+V/h_e$. Its event inequality has no
sample index. An upward coarsening $M\le M'$ is valid only inside the
original moment domain, and the same chosen $M'$ is kept throughout.
The complete high-frequency window is a sum of six pieces. Under the
stated conditions every index-dependent piece decreases for $m\ge N$:
\[
 \sqrt m e^{-cm}\quad(2Nc>1),\qquad 1/m,
 \qquad e^{-km b^2}/\sqrt m.
\]
The comparator has the form
$\sqrt m e^{-m/8}(2r(m/8)+g_c/3)$, where
$r(x)=(x+2)/(x^2+4x+2)>0$ is decreasing and $N\ge5$.
An optional improved piece also has decreasing factor $r(m\ell/4)$,
with $\ell>0$ and its own positive damping and $2Nc>1$ gates.
For the optional zero-piece improvement, put
$f(x)=\sqrt x e^{-x}r(x)$. The sign of its logarithmic derivative is
the sign of
\[
 4-10x-22x^2-13x^3-2x^4,
\]
after multiplication by the positive denominator
$2x(x+2)(x^2+4x+2)$. This polynomial is negative at $x=1/4$ and
strictly decreasing for $x>0$. Thus the improvement is decreasing when
the condition $Nkb^2\ge1/4$ holds. A complete selected window at $N$
is consequently valid later; the minimum of several complete valid
windows is valid without assuming the minimizing window stays the same.
All remaining conditions ($0<a\le1$, $m_e>0$, $0<d_0\le12$,
$1\le M\le111/10$, $0<w\le1/4$, positive sine slope) remain.
These are monotonicity statements for the analytic bounds; numerical
enclosures must bound the complete expressions at the selected floor.

The low-frequency coefficients have index factors
\[
 (m/(m-1))^{3/2},\quad (m/(m-1))^{3/2}/(m-1),\quad
 a_0^{-k},\quad b_0^{-k},\quad m^{-1/2},
\]
where $a_0=(m-1)/(2m)$, $b_0=(45-M)(1-2/m)/96$, and $k>0$.
The first two decrease; $a_0,b_0$ are positive and increase, so their
inverse powers decrease. Thus every sum or product with nonnegative,
index-independent numerators decreases, including the remainder
$e/(2a_0^2)+B^2/(2b_0^3)$ and the additional $L_2/\sqrt m$.
The assertion concerns exactly these analytic envelopes. It requires the
original coefficient theorem, its nonnegative numerators, actual support,
and every paired eligibility gate; it does not repair an invalid paired
rectangle. The paired domain for central panels at floor $40$ is
\[
 0<p_l<p_u<1/2,\quad1\le M<3,\quad g\ge0,\quad g^2<2,
 \quad R>2g,\quad M-1-\gamma_l^2\ge0,\quad0\le\theta<1.
\]
The sign-aware extension in Section~\ref{ms:finite:paired-rectangle}
is a distinct theorem, not an automatic alteration of this old contract.

\subsubsection{Density, coherent score corners, and the charge}
\label{ms:finite:floor-phase}
The retained phase theorem, not a separately asserted bound $H^2\le5r$,
supplies the density enclosure
\[
 \phi(H/\sqrt m)\in[\phi_{0,l}(1-5r/(2N)),2/5]
 \quad(m\ge N)
\]
when its raw and positivity gates hold. This fixed interval contains the
actual interval because its lower endpoint increases with $m$.
An own-score numerator
$r/N-\phi_{0,l}(1-5r/(2N))D_l/6$ bounds its actual counterpart above
for fixed $r,D_l\ge0$: its derivative in $1/N$ is positive.
For $g_0=(79/25)\rho-3-d$ and $t=(d-4/25)_+$, the general phase
curvature has derivative $(g_0+t)/2\ge0$ in the lower density.
For $d\le4/25$, $g_0\ge4/25-d$; for $d\ge4/25$ the sum is
$(79/25)(\rho-1)$. The sign-corrected paired curvature is
\begin{equation}\label{ms:finite:signed-curvature}
 C\rho+\tfrac12\min\{\phi_lg_0,(2/5)g_0\}.
\end{equation}
If $g_0\ge0$ it increases with $\phi_l$; if $g_0<0$ it is constant
at the upper-density endpoint. This proves the floor direction without
calling a negative coefficient positive. Every original positivity gate,
the check against the original valid curvature, both signed linear endpoints,
and the conserved addition with positive coefficient
$(0.4095\rho-d/15)(1-2.5r/N)/2$ remain mandatory.

For a sharp contrary score $C>\CE$ and certified
$\CE<40974/100000$, one has
\[
 m(C-\CE)\rho\ge N\max\{0,C-40974/100000\}\rho_l.
\]
When $C\le40974/100000$ the right side is zero; otherwise each
replacement decreases a nonnegative factor. The breakpoint must be
included in the score cover. On each side the subtracted charge is affine,
so the existing convex score-corner argument remains valid. A positive
$K/n$ target or its saved score floor cannot be transferred by this sharp
$K=0$ argument. The coherent $(C,\rho)$ corners and both signs of $H$
are retained before maximizing. The resulting budget exports
$mD\le Nd$ as in \eqref{ms:finite:weighted-invariant}.
An alternative scalar phase bound
$Q(H)=997Nd/15000+\epsilon(H_{\max}^2-H^2)$, $\epsilon\ge0$,
is at least $997Nd/15000$ on its acquired phase interval; its actual
bound $15000Q(H)/(997m)$ transfers as well. Its artificial curvature
is not a new physical assertion.

\subsubsection{The contracts at floors 25 and 31}
\label{ms:finite:floor25-31}
The floor $25$ construction is restricted to
$20\le n\le40$ and proposed transfer scope $25\le m\le40$;
the intended finite consumers are $m=26,\ldots,30$, with sharp actual
predecessors. Its retained INITIAL variant uses the older gap distance;
a product-gap formula from a different family is not retroactively part of
that arithmetic. A new exact family must establish its own complete arithmetic
before the preceding finite induction can consume it.
The B0 extension adds only the three index-free mean caps above and finite
same-law feedback. It therefore preserves the weighted invariant if its
incoming states and every actual support radius are independently licensed.
Its composition requires all twelve selected support states at each
of the five actual indices, complete initial/resource/local covers, and all
terminal necessary-condition exclusions. If those covers are all empty,
no count forest is needed.

For floor $31$, the eight-pass INITIAL argument uses $N=31$, $K=0$ and
an acquired own-score cap at most $432313/10^6$. The scalar directions in
Sections~\ref{ms:finite:floor-envelopes}--\ref{ms:finite:floor-phase}
prove the INITIAL weighted interpretation for $m\ge31$ on its same
moment/support branch. Each subsequent resource refinement must also
be valid at the actual index.
The intended finite bridge is $31$--$39$, and still requires actual-index
support acquisition, its contact allowance, a valid actual own cap, every
downstream refinement, and complete actual-index counts for surviving rows.
The composition applies on the intersection of all these index domains.

\subsubsection{Support and local comparisons at indices 40, 41 and 42}
\label{ms:finite:floor40}
The following retained contract is only for $m\in\{40,41,42\}$.
It assumes genuine global sharp bounds through 39 (including 38), then uses
successive weak bounds at $40$ and $41$ for the next predecessors. Consequently the
predecessor decrement for the same law is zero. A conservative scalar witness
retains the older allowance $19/100$ and gives
\[
 e=\frac{13/2}{15\cdot38\cdot6}
   +\frac{(2/25)(13/2)}{38^2}=\frac{1469}{649800},
\]
\[
 \begin{gathered}
 U=\frac{10131}{10000}\left(\frac{40974}{100000}+\frac{19}{4000}\right)+e
 =\frac{1371664491931}{3249000000000}<0.4226,\\
 42(U-0.4095)=\frac{288392943517}{541500000000}<0.54.
 \end{gathered}
\]
The inputs $\sqrt m<13/2$, $\sqrt{m-1}>6$ and $m-1\ge38$ hold
on these indices. These computations supply the score and comparison
caps only with the separate predecessor-smoothing theorem. They do not
assume the desired bound at the current index.

The first support bound requires a separate two-step indexed proof at floor $40$,
upper index $42$, zero allowance and radius $3.84$. Every later central-panel
Gaussian support cell instead has the form
\begin{align}
 B_m={}&L_P+\frac{0.54 Q_+}{m}
 +0.4\left(\sum_{j=1}^5\frac{M_j}{m^j}+\frac{T}{m^6}\right)
 +\rho\left(\sqrt{\frac m{m-1}}\delta_m+\frac{0.384\cdot0.4226}{m}\right),
 \label{ms:finite:central-support}
\end{align}
where $Q_+,M_j,T\ge0$, $L_P$ is index-free and may be negative,
and $\delta_m\le0.19$. Holding the same acquired moment, support,
signed-beta and threshold cell bounds fixed, every positive term decreases,
so $B_m\le B_{40}<0$ on each whole closed cell. The prior support
closes the exterior remainder. This argument has discarded favourable
Gaussian terms. It does not justify a primitive retaining such a term;
that is why the indexed initial proof is separate. LEFT has its own
explicit zero-allowance support proof for $40\le m\le42$ at every
selected radius.

For later Sperner confinement and profiles, preserve the complete gates:
$p$ separated from zero and $1/2$, $\theta<1$, positive gaps, support,
buffers, disjointness, positive curvature, both signs and every closed
threshold/noise cell. A fixed deletion factor strictly larger than
$\sqrt{\max((40/39)^3,40/38)}$ bounds both deletions for $m\ge40$.
The floor binomial-mode interval at $j=38$ must contain the actual one at
$j=m-2$, with the original variance and Stirling bounds. This containment
and the analytic mode estimate are explicit prerequisites; no mode is omitted.
Fixed valid derivative bounds $L,U$ remain such bounds: all subtracted
errors have nonnegative numerator divided by $m$, and all added errors
are bounded above by their floor $40$ versions. The lower-ray proof integrates
the same $L$ on both intervals with positive checked width; the upper ray
retains the positive allocation denominator $L+U$.

Let $r_0\ge0$ be the common raw budget before a binary penalty $P\ge0$.
Feasibility gives $r_0-mP\ge0$. For $z,\mathrm{Dist}\ge0$ the
error has the upper bound
\[
 (zr_0+\mathrm{Dist})/m-zP+\mathrm{Bad}\,b_m.
\]
Only the nonnegative inverse numerator is weakened to floor $40$; the term
$-zP$ is retained unchanged. A floor test $r_0-40P\le0$ is safe
with the original strict-margin convention. The bad-pair fraction has
$d\ge0$, $\tau>0$, and
\[
 g_m(d)=\frac d{m\tau+d}\le g_{40}(d)
 \le\frac{40\tau}{(40\tau+s)^2}d+\frac{s^2}{(40\tau+s)^2},\quad s\ge0.
\]
The second inequality is the tangent bound for the concave $g_{40}$,
including $s=0$. The actual feasible moment/noise triangle lies in the
weaker floor triangle, with nonnegative vertex coefficients and
$\mathrm{Bad}$. An independently valid Cantelli bound may be minimized
with this complete bound. The conversion from a resulting amplitude cap to
an $X$-support radius uses the original $p,\ell$ geometry, with no index.
Every feedback step then uses this acquired radius and the weighted invariant.

The final small- and intermediate-noise comparisons keep the normalization entropy segment
$p\pm\delta/40$, which contains $p\pm\delta/m$, and all positive
log-factorial, variance and ratio errors. Normalized mass lies in the
original interval $[b_0,1]$. The signed Gaussian derivatives are bounded
by the expressions $1-A_m/m$ on the positive phase and
$1/(1-\eta_m/m)$ on the negative phase, where
$0\le A_m\le A_{40}$, $0\le\eta_m\le\eta_{40}<40$.
All raw threshold and phase domains and $40v>\delta^2$ are retained.
A cubic correction is either nonnegative or minimized over its complete
positive multiplier interval before entering the repair ceiling.
The full-count inverse correction $I$ may have either sign; exactly
\[
 I/m\le\max\{I/40,I/42\},\qquad40\le m\le42.
\]
For a bracket of either sign and mass in $[b_0,1]$, use
$\max\{b_0\,\mathrm{bracket},\mathrm{bracket}\}$. These two
sign-sensitive maxima retain the restriction $40\le m\le42$. The binomial gap is the minimum of the three actual-index gaps,
not a substituted zero-noise Bernoulli theorem.

The small-noise comparison covers $0<\lambda\le0.05$, and the
intermediate-noise comparison covers $0.05$--$0.08$ at every phase.
On $0.08$--$0.10$, the intermediate-noise comparison covers absolute
phase $0.19$--$0.5$ and the full-count comparison covers its complement.
On $0.10$--$0.12$ their common phase boundary is $0.27$; above $0.12$
the full-count comparison covers all phases. Intersect every interval with the same
row's noise cap, retain both signs of $H$ and phase and all closed seams.
Zero noise is covered by the separate Bernoulli theorem. No smaller-noise
interval is lost at zero.

\begin{condition}[Finite composition at floor 40]\label{ms:finite:floor40-condition}
Supply the actual predecessors, the separate high-moment theorem restricted
to $40,41,42$, all ten acquisition conditions with their
complete hypotheses, the separate indexed initial-support bound, the
complete full-count comparison rows and both binding conditions linking
them to the acquisition bounds, and every small-noise, intermediate-noise,
binomial and coverage condition. Establish all analytic prerequisites and
inequalities used above for the specified bounds and the same maximizing law.
The condition retains all 304 original row-verification obligations. The complete acquisition dispositions
must cover every probability panel outside the left range $[141/500,19/50]$ and
central range $[19/50,119/250]$ as well as all moment panels, including the
shared $\rho=6/5$ boundary.
\end{condition}
Under this condition, finite induction excludes the extremal pair successively at
$40$, $41$ and $42$, without a tail theorem starting at $43$. The six-piece high-frequency window, extended
paired low-frequency bound, signed-variance bound and signed phase estimate
must be the
ones whose inequalities were used. In particular, the inactive base damping
$43/96$ and an unverified paired domain cannot substitute for the required
$(45-M)/96$ domain. The signed-variance bound has no extra sample-index
term, but requires its same-law signed moment/support gates, positive residual
weights and correct root directions. Numerical conditions retain their
original evaluation indices; the analytic transfer above is restricted to
$m=40,41,42$. Every consumed analytic lemma must hold at these indices.
